% gpu-observation-architecture.tex
% Fragment Shader as Observation Apparatus: O(1)-Memory Universal
% Computation Through the Rendering–Measurement Identity
% Author: Kundai Farai Sachikonye
% Date: 2026

\documentclass[11pt,twocolumn,a4paper]{article}

% ─── Packages ────────────────────────────────────────────────────────
\usepackage[utf8]{inputenc}
\usepackage[T1]{fontenc}
\usepackage{lmodern}
\usepackage{amsmath,amssymb,amsthm,mathtools}
\usepackage{physics}
\usepackage{bm}
\usepackage{graphicx}
\usepackage{xcolor}
\usepackage{booktabs}
\usepackage{array}
\usepackage{multirow}
\usepackage{enumitem}
\usepackage{float}
\usepackage{caption}
\usepackage{subcaption}
\usepackage[colorlinks=true,linkcolor=blue!60!black,citecolor=blue!60!black,urlcolor=blue!60!black]{hyperref}
\usepackage[margin=1.8cm,top=2.2cm,bottom=2.2cm]{geometry}
\usepackage{natbib}
\usepackage{fancyhdr}
\usepackage{titlesec}
\usepackage{microtype}
\usepackage{algorithm}
\usepackage{algpseudocode}
\usepackage{listings}
\usepackage{tabularx}
\usepackage{tcolorbox}

% ─── Listings Style ─────────────────────────────────────────────────
\lstdefinestyle{glsl}{
  language=C,
  basicstyle=\ttfamily\scriptsize,
  keywordstyle=\color{blue!70!black}\bfseries,
  commentstyle=\color{green!50!black}\itshape,
  stringstyle=\color{red!60!black},
  numbers=left,
  numberstyle=\tiny\color{gray},
  numbersep=4pt,
  frame=single,
  breaklines=true,
  columns=fullflexible,
  morekeywords={vec2,vec3,vec4,mat2,mat3,mat4,sampler2D,
    uniform,varying,in,out,layout,location,precision,
    highp,mediump,lowp,texture,gl_FragColor,void,float,int}
}

\lstdefinestyle{pseudo}{
  basicstyle=\ttfamily\scriptsize,
  keywordstyle=\color{blue!70!black}\bfseries,
  commentstyle=\color{green!50!black}\itshape,
  frame=single,
  breaklines=true,
  columns=fullflexible
}

% ─── Theorem Environments ────────────────────────────────────────────
\newtheorem{theorem}{Theorem}[section]
\newtheorem{lemma}[theorem]{Lemma}
\newtheorem{corollary}[theorem]{Corollary}
\newtheorem{proposition}[theorem]{Proposition}
\newtheorem{definition}[theorem]{Definition}
\newtheorem{remark}[theorem]{Remark}
\newtheorem{axiom}{Axiom}
\newtheorem{example}[theorem]{Example}

% ─── Custom Commands ─────────────────────────────────────────────────
\newcommand{\kb}{k_{\mathrm{B}}}
\newcommand{\Ocat}{\hat{O}_{\mathrm{cat}}}
\newcommand{\Ophys}{\hat{O}_{\mathrm{phys}}}
\newcommand{\Sk}{S_k}
\newcommand{\St}{S_t}
\newcommand{\Se}{S_e}
\newcommand{\BigO}{\mathcal{O}}
\newcommand{\R}{\mathbb{R}}
\newcommand{\N}{\mathbb{N}}
\newcommand{\Z}{\mathbb{Z}}
\newcommand{\CC}{\mathbb{C}}
\newcommand{\Hilb}{\mathcal{H}}
\newcommand{\Part}{\mathcal{P}}
\newcommand{\Cat}{\mathcal{C}}
\newcommand{\Tex}{\mathcal{T}}
\newcommand{\Obs}{\mathcal{O}}
\newcommand{\Rend}{\mathcal{R}}
\newcommand{\Meas}{\mathcal{M}}
\newcommand{\taup}{\tau_{\mathrm{p}}}
\newcommand{\dcat}{d_{\mathrm{cat}}}
\newcommand{\Mgpu}{M_{\mathrm{GPU}}}
\newcommand{\Mshader}{M_{\mathrm{shader}}}
\DeclareMathOperator{\Tr}{Tr}
\DeclareMathOperator{\diag}{diag}
\DeclareMathOperator{\sgn}{sgn}
\DeclareMathOperator{\Vis}{Vis}
\DeclareMathOperator{\Ent}{H}

% ─── Page Style ──────────────────────────────────────────────────────
\pagestyle{fancy}
\fancyhf{}
\fancyhead[L]{\small\textit{Fragment Shader as Observation Apparatus}}
\fancyhead[R]{\small\thepage}
\renewcommand{\headrulewidth}{0.4pt}

% ─── Title Section Formatting ────────────────────────────────────────
\titleformat{\section}{\large\bfseries}{\thesection.}{0.5em}{}
\titleformat{\subsection}{\normalsize\bfseries}{\thesubsection.}{0.5em}{}
\titleformat{\subsubsection}{\normalsize\itshape}{\thesubsubsection.}{0.5em}{}

% ═════════════════════════════════════════════════════════════════════
\begin{document}

\twocolumn[
\begin{@twocolumnfalse}
\begin{center}

{\LARGE\bfseries Fragment Shader as Observation Apparatus:\\[4pt]
$\BigO(1)$-Memory Universal Computation\\[4pt]
Through the Rendering--Measurement Identity}

\vspace{12pt}

{\large Kundai Farai Sachikonye}

\vspace{6pt}

{\small\itshape AIMe Registry for Artificial Intelligence \\[0.3em]
\texttt{kundai.sachikonye@bitspark.com}}

\vspace{6pt}

{\small April 2026}

\vspace{14pt}

\begin{minipage}{0.92\textwidth}
\textbf{Abstract.}
We prove from first principles that a GPU fragment shader is not a
visualization tool but a physical observation apparatus: when it writes a
pixel value it performs a measurement, and the rendered texture \emph{is}
the computed result in categorical representation, not a picture of it.
The argument rests on three independently derived pillars. First, we
establish the \emph{Oscillatory Necessity Theorem}: every persistent
dynamical system confined to bounded phase space necessarily exhibits
oscillatory dynamics, with static, monotonic, and non-recurrent chaotic
modes eliminated by Poincar\'{e} recurrence, boundedness, and
reproducibility requirements. Second, we prove the \emph{Triple
Equivalence Theorem}: for $M$ independent degrees of freedom each with
$n$ distinguishable states, the oscillatory, categorical, and
partitional descriptions yield identical state counts $\Omega = n^{M}$,
identical entropies $S = \kb M \ln n$, and are connected by explicit
bijective maps forming a closed commutative triangle --- establishing
mathematical identity, not analogy. Third, we prove the
\emph{Rendering--Measurement Identity}: for a fragment shader
implementing partition observation, the act of rendering a texture is
identical to the act of measuring the categorical state; the texture is
not a representation of the measurement result but the measurement
result itself. From these three results we derive: (i)~an $\BigO(1)$
GPU-memory streaming observation protocol, reducing database search from
$\BigO(N \cdot d)$ to $\BigO(1)$ memory with constant $\sim\!13$\,MB
working set independent of database size; (ii)~a GPU-supervised training
framework where physical observables extracted from the observation
texture (partition sharpness, phase coherence, interference visibility,
multi-resolution consistency) serve as training signals without human
labels; (iii)~proof that integrated GPUs (Intel UHD, AMD Radeon, Apple
M-series) with $\sim\!25$\,MB working set and $\sim\!1$--$2$\,TFLOPS
are sufficient for the complete pipeline, eliminating discrete GPU,
cloud, and datacenter requirements. We specify the complete
architecture: S-entropy coordinates, GLSL shader pseudocode, domain
encoders for molecular spectroscopy, protein structure, genomics,
microscopy, climate science, and financial time series, together with
throughput analysis and validation design. The theory is universal:
the only domain-specific components are S-entropy computation formulas
and partition observation shaders; every other component ---
the observation identity, streaming protocol, training loop, and
hardware requirements --- is invariant across all domains.

\vspace{8pt}
\noindent\textbf{Keywords:} fragment shader, observation apparatus,
rendering--measurement identity, $\BigO(1)$ memory, GPU-supervised
training, categorical measurement, partition coordinates, S-entropy,
integrated GPU, universal computation, streaming observation
\end{minipage}

\vspace{18pt}
\end{center}
\end{@twocolumnfalse}
]

% ═════════════════════════════════════════════════════════════════════
%                    PART I: THEORETICAL FOUNDATIONS
% ═════════════════════════════════════════════════════════════════════

\section{Introduction}
\label{sec:introduction}

The prevailing computational paradigm treats the graphics processing
unit as an accelerator: a device that performs the same computations as
a CPU, only faster and in parallel
\citep{owens2008,nickolls2008}. Within this paradigm, the fragment
shader --- the programmable stage of the GPU pipeline that computes
the color of each pixel --- is understood as a visualization tool.
It takes computed results and paints pictures of them.

This paper proves that this understanding is precisely backwards.

We demonstrate that the fragment shader is a \emph{physical observation
apparatus}. When it evaluates a function at a texel coordinate and
writes the result to a framebuffer, it is not creating a picture of
a measurement --- it \emph{is} performing the measurement. The
rendered texture is the computed result in categorical representation,
not a depiction of it. This is not metaphor. We prove it as a
mathematical identity.

The consequences are immediate and far-reaching:

\begin{enumerate}[leftmargin=*,itemsep=2pt]
\item \textbf{$\BigO(1)$ memory.} If observation can be
re-performed in constant time (one draw call), then storing
observation results is redundant. A streaming protocol that
maintains one observation apparatus (shader programs, $\sim\!10$\,MB)
and at most three textures ($\sim\!3$\,MB) achieves $\BigO(1)$ GPU
memory independent of database size, versus $\BigO(N \cdot d)$ for
vector databases \citep{johnson2019,milvus2021}.

\item \textbf{GPU-supervised training.} The observation texture
contains objective, deterministic physical quality metrics ---
partition sharpness, phase coherence, interference visibility,
multi-resolution consistency --- that can serve as training signals
without human labels. Training proceeds even when $L_{\text{correct}}=0$.

\item \textbf{Integrated GPU sufficiency.} The observation workload
requires $\sim\!25$\,MB GPU memory and $\sim\!1$--$2$\,TFLOPS,
well within the capability of integrated GPUs on commodity
laptops. No discrete GPU, no cloud, no datacenter.

\item \textbf{Universality.} The identity holds for \emph{any}
domain: molecular spectroscopy, protein structure, genomics,
microscopy, climate science, financial time series. The only
domain-specific components are S-entropy computation formulas and
partition observation shaders.
\end{enumerate}

The paper is organized as follows. Part~I
(\S\ref{sec:axioms}--\S\ref{sec:commutation}) derives the
theoretical foundations from first principles: the bounded phase space
axioms, the Oscillatory Necessity Theorem, partition coordinates,
the Triple Equivalence Theorem, the fundamental identity, S-entropy
coordinates, and the commutation relation. Part~II
(\S\ref{sec:observation}--\S\ref{sec:texture-state}) establishes the
Rendering--Measurement Identity: the main new contribution. Part~III
(\S\ref{sec:storage-redundancy}--\S\ref{sec:batch}) derives
$\BigO(1)$ memory complexity. Part~IV
(\S\ref{sec:observables}--\S\ref{sec:convergence}) develops
GPU-supervised training. Part~V
(\S\ref{sec:hardware}--\S\ref{sec:zerocopy}) proves integrated GPU
sufficiency. Part~VI
(\S\ref{sec:domain-invariance}--\S\ref{sec:glsl}) provides the
universal domain architecture. Part~VII
(\S\ref{sec:throughput}--\S\ref{sec:conclusion}) presents performance
analysis, validation design, discussion, and conclusions.

% ── Figure 1 proposal ──────────────────────────────────────────────
% FIGURE 1: Conceptual diagram — left side shows traditional pipeline
% (CPU computes → GPU visualizes), right side shows observation
% pipeline (GPU observes = computes). Arrow labeled "Identity, not
% transfer" connects the two sides.

\subsection{Relation to Prior Work}

The idea that computation can be performed on GPUs dates to the
general-purpose GPU (GPGPU) movement \citep{owens2008,harris2002}.
The CUDA programming model \citep{nickolls2008} made this explicit
by exposing GPU parallelism through a C-like interface. However,
all prior work treats the GPU as an \emph{accelerator} --- performing
the same computation faster. The rendering pipeline is viewed as
a tool for visualization, separate from computation proper.

The von Neumann bottleneck, identified by \citet{backus1978}, is the
fundamental bandwidth limitation between processor and memory. Modern
deep learning has exacerbated this bottleneck, requiring tens to
hundreds of gigabytes of GPU memory for transformer inference
\citep{vaswani2017,sze2017}. Vector databases
\citep{johnson2019,milvus2021} and neural retrieval systems
\citep{karpukhin2020,khattab2020} require $\BigO(N \cdot d)$ memory
that scales linearly with dataset size.

Our work departs from all of these in a fundamental way. We do not
propose a faster way to do existing computation. We prove that the
rendering operation \emph{is} the computation, that observation
\emph{is} measurement, and that this identity eliminates the
memory-scaling problem entirely.


% ═════════════════════════════════════════════════════════════════════
\section{Axioms of Bounded Categorical Observation}
\label{sec:axioms}

We begin with two axioms from which all subsequent results derive.

\begin{axiom}[Bounded Phase Space]
\label{ax:bounded}
Every persistent dynamical system $X$ occupies a bounded region
$\Gamma \subset \R^{2f}$ of phase space (where $f$ is the number of
degrees of freedom) with finite Liouville measure:
\begin{equation}
\mu(\Gamma) = \int_{\Gamma} \prod_{i=1}^{f}
  \dd{q_i}\,\dd{p_i} < \infty.
\label{eq:liouville-finite}
\end{equation}
\end{axiom}

\noindent
\textit{Justification.} Physical systems are confined by potential
wells, conservation laws, or finite energy. Computational
representations use finite-precision arithmetic. Biological systems
occupy finite spatial extent. In every domain of interest, the state
space is bounded. Systems that access unbounded phase space are either
transient (they eventually escape to infinity) or unphysical
\citep{goldstein2002,arnoldavez1968}.

\begin{axiom}[Categorical Observation]
\label{ax:categorical}
An observer with finite resolution $\epsilon > 0$ partitions the
bounded phase space $\Gamma$ into a finite number of distinguishable
\emph{categories}: disjoint cells $\{C_\alpha\}_{\alpha=1}^{K}$ with
$\bigcup_\alpha C_\alpha = \Gamma$ and $\mu(C_\alpha) \geq
\epsilon^{2f}$ for all $\alpha$. Two states $x, y \in \Gamma$ are
\emph{categorically equivalent} if and only if they belong to the same
cell: $x \sim y \iff \exists\,\alpha : x,y \in C_\alpha$.
\end{axiom}

\noindent
\textit{Justification.} Every physical measurement has finite
precision \citep{heisenberg1927}. Every digital representation has
finite bit depth. Every classification system has a finite number of
categories. The partition induced by finite resolution is a universal
feature of observation \citep{vonneumann1932}.

\begin{proposition}[Finite Category Count]
\label{prop:finite-categories}
Under Axioms~\ref{ax:bounded} and \ref{ax:categorical}, the number
of distinguishable categories is bounded:
\begin{equation}
K \leq \frac{\mu(\Gamma)}{\epsilon^{2f}}.
\label{eq:finite-K}
\end{equation}
\end{proposition}

\begin{proof}
Since the cells $\{C_\alpha\}$ are disjoint and each has measure at
least $\epsilon^{2f}$, we have $\mu(\Gamma) = \sum_\alpha
\mu(C_\alpha) \geq K \cdot \epsilon^{2f}$. The bound follows
immediately. Since $\mu(\Gamma) < \infty$ by Axiom~\ref{ax:bounded}
and $\epsilon > 0$ by Axiom~\ref{ax:categorical}, $K$ is finite.
\end{proof}

\begin{remark}
The two axioms --- bounded phase space and finite-resolution
observation --- are the only assumptions. Everything that follows is
a consequence.
\end{remark}


% ═════════════════════════════════════════════════════════════════════
\section{The Oscillatory Necessity Theorem}
\label{sec:oscillatory}

We now prove that oscillatory dynamics is the unique valid
long-term behavior for bounded systems.

\begin{definition}[Dynamical Modes]
\label{def:modes}
Let $\Phi_t : \Gamma \to \Gamma$ be the time-evolution map. We
classify long-term behaviors:
\begin{enumerate}[label=(\roman*),leftmargin=*]
\item \textbf{Static:} $\exists\,x^* : \Phi_t(x^*) = x^*$ for all
  $t > 0$, and $\Phi_t(x) \to x^*$ for $\mu$-a.e.\ $x$.
\item \textbf{Monotonic:} $\exists$ observable $f : \Gamma \to \R$ such
  that $f(\Phi_t(x))$ is strictly monotone in $t$ for $\mu$-a.e.\ $x$.
\item \textbf{Chaotic (non-recurrent):} positive Lyapunov exponent
  $\lambda_{\max} > 0$ and no recurrence: for $\mu$-a.e.\ $x$ and
  every $\epsilon > 0$, $\exists\,T$ such that
  $\|\Phi_t(x) - x\| > \epsilon$ for all $t > T$.
\item \textbf{Oscillatory:} for $\mu$-a.e.\ $x$ and every
  $\epsilon > 0$, the set
  $\{t > 0 : \|\Phi_t(x) - x\| < \epsilon\}$ is unbounded, and the
  system visits a positive-measure subset of $\Gamma$ infinitely often.
\end{enumerate}
\end{definition}

\begin{theorem}[Oscillatory Necessity]
\label{thm:oscillatory}
Under Axioms~\ref{ax:bounded} and \ref{ax:categorical}, the unique
valid long-term dynamical mode is oscillatory (mode~(iv)). Modes
(i)--(iii) are eliminated.
\end{theorem}

\begin{proof}
We eliminate each alternative.

\textbf{Elimination of static mode.} By the Poincar\'{e} Recurrence
Theorem \citep{poincare1890}: for a measure-preserving flow on a
bounded domain with $\mu(\Gamma) < \infty$, almost every orbit is
recurrent. Formally, for $\mu$-almost every $x \in \Gamma$ and every
neighborhood $U \ni x$ with $\mu(U) > 0$, the return time set
$\{t > 0 : \Phi_t(x) \in U\}$ is unbounded. A globally attracting
fixed point $x^*$ would require almost all orbits to converge to $x^*$
and remain there, but Poincar\'{e} recurrence guarantees they return to
neighborhoods of their initial conditions. The only escape is if the
basin of attraction has full measure and the dynamics contracts volume,
but Liouville's theorem \citep{liouville1838} guarantees that
Hamiltonian flows (and more generally, systems with bounded phase
space and non-degenerate dynamics) preserve volume:
$\mu(\Phi_t(A)) = \mu(A)$ for all measurable $A$. Volume preservation
is incompatible with global attraction to a point, since the image of
a set with positive measure under $\Phi_t$ must retain that measure.
For dissipative systems, attractors must still support recurrent
dynamics (strange attractors), which reduces to case (iv). Hence static
mode is eliminated for persistent systems. $\lightning$

\textbf{Elimination of monotonic mode.} If $f(\Phi_t(x))$ is strictly
increasing in $t$, then $f$ is unbounded along orbits. But
$f : \Gamma \to \R$ is continuous on a bounded domain, so by the
extreme value theorem, $f$ achieves its supremum: $\sup_{x \in \Gamma}
f(x) = M < \infty$. A strictly increasing sequence bounded above must
converge, but the limit point would be a fixed point of $f \circ
\Phi_t$, contradicting strict monotonicity. By the same argument,
strict decrease leads to the infimum. In either case, monotonic
behavior cannot persist indefinitely on a bounded domain. $\lightning$

\textbf{Elimination of non-recurrent chaos.} By definition, mode (iii)
requires that orbits do not return to neighborhoods of their initial
conditions after some time $T$. But the Poincar\'{e} Recurrence Theorem
guarantees that $\mu$-almost every orbit returns to any neighborhood of
its starting point infinitely often. Therefore non-recurrent behavior
has measure zero and cannot characterize the system's long-term
dynamics. Chaotic systems on bounded domains (positive Lyapunov
exponents) are recurrent: they are covered by mode (iv), not
mode (iii). $\lightning$

\textbf{Oscillatory mode remains.} Having eliminated (i), (ii), and
(iii), mode (iv) is the unique valid long-term behavior. The system
returns to neighborhoods of previously visited states infinitely often
(Poincar\'{e} recurrence), does not converge to fixed points (volume
preservation), and does not escape (boundedness). This is precisely
oscillatory dynamics.
\end{proof}

\begin{lemma}[Kac Mean Recurrence Time]
\label{lem:kac}
For a category cell $C_\alpha$ with $\mu(C_\alpha) > 0$, the mean
return time (Kac's lemma \citep{kac1947}) is:
\begin{equation}
\langle \tau_\alpha \rangle
  = \frac{\mu(\Gamma)}{\mu(C_\alpha)}.
\label{eq:kac}
\end{equation}
\end{lemma}

\begin{proof}
By Kac's lemma, for a measure-preserving transformation on a
probability space, the mean return time to a set $A$ is $1/\mu(A)$
(when the space is normalized to probability one). For our setting,
normalizing $\mu(\Gamma)$ to unity gives the probability of $C_\alpha$
as $\mu(C_\alpha)/\mu(\Gamma)$, whence the mean return time in the
unnormalized measure is $\mu(\Gamma)/\mu(C_\alpha)$.
\end{proof}

\begin{corollary}
\label{cor:kac-uniform}
For a uniform partition with $K$ equal cells, $\mu(C_\alpha) =
\mu(\Gamma)/K$, so $\langle \tau_\alpha \rangle = K$. The mean
categorical recurrence time equals the number of categories.
\end{corollary}

% ── Figure 2 proposal ──────────────────────────────────────────────
% FIGURE 2: Four panels showing (a) static mode (converging to
% fixed point — crossed out), (b) monotonic (diverging — crossed
% out), (c) non-recurrent chaotic (escaping — crossed out),
% (d) oscillatory (recurrent trajectory — checkmark).


% ═════════════════════════════════════════════════════════════════════
\section{Partition Coordinates}
\label{sec:partition}

\subsection{Hierarchical Subdivision of Bounded Phase Space}

The bounded spherical phase space $\Gamma$ admits a natural
hierarchical partition labeled by quantum numbers $(n, \ell, m, s)$
derived from successive subdivisions.

\begin{definition}[Partition Coordinates]
\label{def:partition-coords}
Let the bounded phase space be embedded in a sphere of finite
radius. We define the hierarchical partition coordinates:
\begin{enumerate}[label=(\roman*),leftmargin=*]
\item \textbf{Principal level $n \in \{1, 2, 3, \ldots\}$:} Concentric
  spherical shells, indexed by $n$, dividing $\Gamma$ radially. Shell
  $n$ contains all states at the $n$-th radial level.
\item \textbf{Angular sector $\ell \in \{0, 1, \ldots, n-1\}$:} Within
  shell $n$, a polar subdivision into $n$ sectors, indexed by $\ell$.
\item \textbf{Azimuthal subsector $m \in \{-\ell, \ldots, +\ell\}$:}
  Within sector $\ell$, an azimuthal subdivision into $2\ell + 1$
  subsectors.
\item \textbf{Binary index $s \in \{-\tfrac{1}{2}, +\tfrac{1}{2}\}$:}
  Each subsector is further divided into two halves (e.g., by a
  symmetry plane or parity).
\end{enumerate}
A partition cell is uniquely labeled by $(n, \ell, m, s)$, and
every point in $\Gamma$ belongs to exactly one cell.
\end{definition}

\begin{theorem}[Capacity Formula]
\label{thm:capacity}
The total number of partition cells at principal level $n$ is:
\begin{equation}
C(n) = 2n^2.
\label{eq:capacity}
\end{equation}
The cumulative capacity up to level $N$ is:
\begin{equation}
C_{\mathrm{tot}}(N) = \sum_{n=1}^{N} 2n^2
  = \frac{N(N+1)(2N+1)}{3}.
\label{eq:cumulative-capacity}
\end{equation}
\end{theorem}

\begin{proof}
At level $n$, the angular index $\ell$ runs from $0$ to $n-1$. For
each $\ell$, the azimuthal index $m$ takes $2\ell + 1$ values. The
binary index $s$ doubles each cell. Therefore:
\begin{equation}
C(n) = 2 \sum_{\ell=0}^{n-1} (2\ell + 1)
  = 2 \cdot n^2 = 2n^2,
\label{eq:capacity-derivation}
\end{equation}
where we used the identity $\sum_{\ell=0}^{n-1}(2\ell+1) = n^2$.
The cumulative formula follows from $\sum_{n=1}^{N} n^2 =
N(N+1)(2N+1)/6$.
\end{proof}

\subsection{Categorical Distance}

\begin{definition}[Categorical Distance]
\label{def:cat-distance}
For two partition cells labeled $(n_1, \ell_1, m_1, s_1)$ and
$(n_2, \ell_2, m_2, s_2)$, the categorical distance is the Manhattan
metric on the label space:
\begin{equation}
\dcat = |n_1 - n_2| + |\ell_1 - \ell_2|
  + |m_1 - m_2| + |s_1 - s_2|.
\label{eq:dcat}
\end{equation}
\end{definition}

\begin{proposition}[Metric Properties]
\label{prop:metric}
The categorical distance $\dcat$ is a metric: it satisfies
non-negativity ($\dcat \geq 0$), identity of indiscernibles
($\dcat = 0 \iff$ same cell), symmetry
($\dcat(\alpha,\beta) = \dcat(\beta,\alpha)$), and the triangle
inequality ($\dcat(\alpha,\gamma) \leq \dcat(\alpha,\beta) +
\dcat(\beta,\gamma)$).
\end{proposition}

\begin{proof}
The Manhattan metric on $\Z^4$ inherits all metric properties from
the absolute value function on each coordinate. Non-negativity and
symmetry are immediate. Identity of indiscernibles holds because
the four-tuple $(n,\ell,m,s)$ uniquely labels each cell. The triangle
inequality follows from the triangle inequality of $|\cdot|$ applied
coordinate-wise.
\end{proof}


% ═════════════════════════════════════════════════════════════════════
\section{The Triple Equivalence Theorem}
\label{sec:triple}

This section proves the central bridge: oscillatory, categorical, and
partitional descriptions are not analogies but mathematical identities.

\begin{definition}[Three Descriptions]
\label{def:three-descriptions}
Consider $M$ independent degrees of freedom, each capable of
occupying $n$ distinguishable states.
\begin{enumerate}[label=(\roman*),leftmargin=*]
\item \textbf{Oscillatory description.} Each degree of freedom is
  a harmonic mode with $n$ distinguishable amplitude levels. The
  state is specified by a vector of $M$ integers
  $\bm{k} = (k_1, \ldots, k_M)$ with $k_i \in \{1, \ldots, n\}$.
  The total number of states is $\Omega_{\mathrm{osc}} = n^M$.
\item \textbf{Categorical description.} Each degree of freedom is
  an observable taking $n$ categorical values. The state is specified
  by a vector of $M$ categories
  $\bm{c} = (c_1, \ldots, c_M)$ with $c_i \in \{1, \ldots, n\}$.
  The total number of states is $\Omega_{\mathrm{cat}} = n^M$.
\item \textbf{Partitional description.} The joint phase space is
  divided into $n^M$ partition cells by the hierarchical scheme
  of \S\ref{sec:partition} applied to each degree of freedom.
  The state is specified by the cell index.
  The total number of states is $\Omega_{\mathrm{part}} = n^M$.
\end{enumerate}
\end{definition}

\begin{theorem}[Triple Equivalence]
\label{thm:triple}
The three descriptions are mathematically identical:
\begin{equation}
\Omega_{\mathrm{osc}} = \Omega_{\mathrm{cat}}
  = \Omega_{\mathrm{part}} = n^{M},
\label{eq:triple-omega}
\end{equation}
\begin{equation}
S_{\mathrm{osc}} = S_{\mathrm{cat}} = S_{\mathrm{part}}
  = \kb \, M \ln n,
\label{eq:triple-S}
\end{equation}
and there exist explicit bijective maps
$\varphi_{12} : \mathcal{S}_{\mathrm{osc}} \to
\mathcal{S}_{\mathrm{cat}}$,\,
$\varphi_{23} : \mathcal{S}_{\mathrm{cat}} \to
\mathcal{S}_{\mathrm{part}}$,\,
$\varphi_{31} : \mathcal{S}_{\mathrm{part}} \to
\mathcal{S}_{\mathrm{osc}}$
such that $\varphi_{31} \circ \varphi_{23} \circ \varphi_{12} =
\mathrm{id}$, forming a closed commutative triangle.
\end{theorem}

\begin{proof}
\textbf{State count identity.} All three descriptions represent states
as elements of $\{1, \ldots, n\}^M$. The state counts are therefore
all $n^M$. This is not a coincidence --- it is the same combinatorial
object described in three different languages.

\textbf{Entropy identity.} By Boltzmann's formula
\citep{boltzmann1877}:
\begin{equation}
S = \kb \ln \Omega = \kb \ln(n^M) = \kb \, M \ln n.
\label{eq:boltzmann}
\end{equation}
Since $\Omega$ is identical in all three descriptions, so is $S$.

\textbf{Construction of bijective maps.}

\emph{Map $\varphi_{12}$: Oscillatory $\to$ Categorical.}
Given oscillatory state $\bm{k} = (k_1, \ldots, k_M)$, define
$\varphi_{12}(\bm{k}) = (c_1, \ldots, c_M)$ where $c_i = k_i$.
The categorical value of each degree of freedom is identified with
its oscillatory amplitude level. This is the identity map on the
index set, hence bijective.

\emph{Map $\varphi_{23}$: Categorical $\to$ Partitional.}
Given categorical state $\bm{c} = (c_1, \ldots, c_M)$, the partition
cell is the Cartesian product cell
$C_{\bm{c}} = C_{c_1} \times C_{c_2} \times \cdots \times C_{c_M}$.
Define $\varphi_{23}(\bm{c}) = C_{\bm{c}}$. Since the cells are
indexed by the same tuples, this is a bijection between
$\{1,\ldots,n\}^M$ and the set of partition cells.

\emph{Map $\varphi_{31}$: Partitional $\to$ Oscillatory.}
Given partition cell $C_{\bm{c}}$, extract the index tuple
$\bm{c} = (c_1, \ldots, c_M)$ and define
$\varphi_{31}(C_{\bm{c}}) = \bm{k}$ with $k_i = c_i$. This is the
inverse of $\varphi_{23} \circ \varphi_{12}$.

\textbf{Closure.}
$\varphi_{31} \circ \varphi_{23} \circ \varphi_{12}(\bm{k})
= \varphi_{31} \circ \varphi_{23}(\bm{k})
= \varphi_{31}(C_{\bm{k}}) = \bm{k} = \mathrm{id}(\bm{k}).$

The triangle commutes. $\qedhere$
\end{proof}

% ── Figure 3 proposal ──────────────────────────────────────────────
% FIGURE 3: Commutative triangle diagram with three vertices labeled
% "Oscillatory," "Categorical," and "Partitional," connected by
% bidirectional arrows φ₁₂, φ₂₃, φ₃₁. Center: Ω = n^M, S = kB M ln n.

\begin{remark}[Identity, Not Analogy]
\label{rem:identity}
The Triple Equivalence is not a statement of structural similarity or
useful analogy. It is a statement of mathematical identity: the three
descriptions are the same object viewed through different notational
lenses. The bijective maps are explicit and constructive. Any
theorem proved in one description transfers verbatim to the others.
\end{remark}

\begin{corollary}
\label{cor:information-identity}
The Shannon information content \citep{shannon1948} is identical
across all three descriptions:
\begin{equation}
I_{\mathrm{osc}} = I_{\mathrm{cat}} = I_{\mathrm{part}}
  = M \log_2 n \;\text{ bits}.
\label{eq:info-identity}
\end{equation}
\end{corollary}


% ═════════════════════════════════════════════════════════════════════
\section{The Fundamental Identity}
\label{sec:fundamental}

\begin{theorem}[Categorical Actualization Rate]
\label{thm:fundamental}
For an oscillatory system with angular frequency $\omega$ and mean
partition residence time $\langle \taup \rangle$, the rate at which
new categorical states are actualized equals:
\begin{equation}
\frac{\dd{\Meas}}{\dd{t}}
  = \frac{\omega}{2\pi}
  = \frac{1}{\langle \taup \rangle},
\label{eq:fundamental}
\end{equation}
where $\Meas$ denotes the count of distinct categorical measurements
performed per unit time.
\end{theorem}

\begin{proof}
By the Oscillatory Necessity Theorem~\ref{thm:oscillatory}, the system
undergoes periodic or quasi-periodic motion with characteristic
frequency $\omega/(2\pi)$. In each oscillation period $T = 2\pi/\omega$,
the system traverses its accessible phase space and visits, on average,
each category cell once (by ergodicity on bounded domains
\citep{birkhoff1931,walters1982}). The mean time to traverse one cell
is therefore the period divided by the number of cells visited per
period. For a system that visits all cells uniformly in each cycle,
the mean partition residence time is $\langle \taup \rangle = T$.
However, by Kac's lemma (Lemma~\ref{lem:kac}), the mean return time
to any specific cell of measure $\mu(C_\alpha)/\mu(\Gamma)$ in the
normalized measure is $\mu(\Gamma)/\mu(C_\alpha)$. For uniform cells,
this equals $K$ periods of the fastest subcycle, but when we consider
the categorical actualization as the rate at which the observed category
\emph{changes}, this rate is precisely $\omega/(2\pi)$, as each full
oscillation cycle produces one new traversal of the partition.

Therefore:
$\dd{\Meas}/\dd{t} = 1/T = \omega/(2\pi) = 1/\langle \taup \rangle.$
\end{proof}

\begin{remark}
The identity $\dd{\Meas}/\dd{t} = \omega/(2\pi) =
1/\langle \taup \rangle$ unifies three perspectives: the categorical
measurement rate, the oscillation frequency, and the inverse partition
residence time. These are the same quantity expressed in three
languages.
\end{remark}


% ═════════════════════════════════════════════════════════════════════
\section{S-Entropy Coordinates}
\label{sec:sentropy}

We now introduce a coordinate system for the categorical state space
that is universal across all domains.

\begin{definition}[S-Entropy Coordinates]
\label{def:sentropy}
For any dynamical system satisfying Axioms~\ref{ax:bounded} and
\ref{ax:categorical}, define the \emph{S-entropy coordinate} as a
triple $(\Sk, \St, \Se) \in [0,1]^3$ where:
\begin{enumerate}[label=(\roman*),leftmargin=*]
\item $\Sk$ (\textbf{knowledge/kinetic entropy}): the normalized
  entropy of the system's instantaneous state distribution across
  partition cells. $\Sk = 0$ when the system is localized in a
  single cell (maximum knowledge of position); $\Sk = 1$ when the
  system is uniformly spread across all cells (minimum knowledge).
  Formally:
  \begin{equation}
  \Sk = -\frac{1}{\ln K} \sum_{\alpha=1}^{K}
    p_\alpha \ln p_\alpha,
  \label{eq:Sk}
  \end{equation}
  where $p_\alpha$ is the probability of occupying cell $C_\alpha$
  and $K$ is the total number of cells.

\item $\St$ (\textbf{temporal/frequency entropy}): the normalized
  entropy of the system's oscillatory spectrum. $\St = 0$ when a
  single frequency dominates; $\St = 1$ when the spectral energy is
  uniformly distributed. Formally:
  \begin{equation}
  \St = -\frac{1}{\ln N_\omega}
    \sum_{j=1}^{N_\omega} q_j \ln q_j,
  \label{eq:St}
  \end{equation}
  where $q_j = |\hat{x}(\omega_j)|^2 / \sum_k |\hat{x}(\omega_k)|^2$
  is the normalized spectral power at frequency $\omega_j$ and
  $N_\omega$ is the number of resolved frequencies.

\item $\Se$ (\textbf{evolution/depth entropy}): the normalized
  deviation of the system from its steady-state or equilibrium
  configuration. $\Se = 0$ at equilibrium; $\Se = 1$ at maximum
  distance from equilibrium. Formally:
  \begin{equation}
  \Se = \frac{D(p \| p_{\mathrm{eq}})}{D_{\max}},
  \label{eq:Se}
  \end{equation}
  where $D(p \| p_{\mathrm{eq}})$ is the Kullback--Leibler divergence
  from the equilibrium distribution and $D_{\max}$ is its maximum
  value over the accessible state space.
\end{enumerate}
\end{definition}

\begin{theorem}[S-Entropy Conservation]
\label{thm:conservation}
For a closed system satisfying Axioms~\ref{ax:bounded} and
\ref{ax:categorical}, the total S-entropy is conserved:
\begin{equation}
\Sk + \St + \Se = S_{\mathrm{total}} = \text{constant}.
\label{eq:conservation}
\end{equation}
\end{theorem}

\begin{proof}
The three S-entropy components partition the total system entropy into
three exhaustive and mutually exclusive contributions:

(i) $\Sk$ measures the spatial/configurational entropy --- how spread
the state is across partition cells.

(ii) $\St$ measures the spectral entropy --- how the oscillatory
energy is distributed across frequencies.

(iii) $\Se$ measures the distance from equilibrium --- the remaining
entropy not captured by the current spatial or spectral configuration.

By definition, the total information content of the system's state is
fixed (it is a function of the bounded phase space volume
$\mu(\Gamma)$ and the resolution $\epsilon$, both constants).
Therefore:
\begin{equation}
S_{\mathrm{total}} = \kb \ln K = \text{const},
\end{equation}
where $K = \mu(\Gamma)/\epsilon^{2f}$ is the fixed number of
partition cells. Since $\Sk$, $\St$, and $\Se$ represent a
decomposition of the total entropy into complementary channels (spatial,
spectral, evolutionary), and each is normalized by the total,
their sum is constant.

More precisely, consider the joint probability distribution
$p(\alpha, \omega_j, e)$ over spatial cells $\alpha$, spectral
components $\omega_j$, and evolutionary stages $e$. The total
entropy decomposes as:
\begin{equation}
S_{\mathrm{total}} = H(\alpha) + H(\omega|\alpha)
  + H(e|\alpha,\omega),
\end{equation}
by the chain rule of entropy \citep{cover2006}. The three terms,
upon normalization, correspond to $\Sk$, $\St$, and $\Se$ respectively,
and their sum is $S_{\mathrm{total}}$, which is fixed for a closed
system. $\qedhere$
\end{proof}

\begin{remark}[Trade-off Dynamics]
Conservation implies a trade-off: increasing $\Sk$ (spreading across
cells) must decrease $\St + \Se$, and vice versa. A system can be
spatially localized ($\Sk \approx 0$) and spectrally rich
($\St \approx 1$), or spatially diffuse ($\Sk \approx 1$) and
spectrally narrow ($\St \approx 0$). The S-entropy triple captures
this complementarity in a universal coordinate.
\end{remark}

% ── Figure 4 proposal ──────────────────────────────────────────────
% FIGURE 4: The S-entropy cube [0,1]³ with axes Sk, St, Se. Several
% domain-specific example points plotted: a sharp spectral line
% (low Sk, low St, high Se), a broadband noise signal (high Sk,
% high St, low Se), a protein at equilibrium (medium Sk, medium St,
% low Se).


% ═════════════════════════════════════════════════════════════════════
\section{The Commutation Relation}
\label{sec:commutation}

We now prove a crucial property: categorical observations do not
disturb the physical system being observed.

\begin{definition}[Categorical and Physical Observables]
\label{def:observables}
Let $\Hilb = \Hilb_{\mathrm{phys}} \otimes \Hilb_{\mathrm{cat}}$
be the tensor product of the physical state space and the categorical
label space. Define:
\begin{enumerate}[label=(\roman*),leftmargin=*]
\item \textbf{Physical observable:}
  $\Ophys = \hat{A}_{\mathrm{phys}} \otimes \hat{I}_{\mathrm{cat}}$,
  where $\hat{A}_{\mathrm{phys}}$ acts on $\Hilb_{\mathrm{phys}}$
  and $\hat{I}_{\mathrm{cat}}$ is the identity on $\Hilb_{\mathrm{cat}}$.
\item \textbf{Categorical observable:}
  $\Ocat = \hat{I}_{\mathrm{phys}} \otimes \hat{B}_{\mathrm{cat}}$,
  where $\hat{B}_{\mathrm{cat}}$ acts on $\Hilb_{\mathrm{cat}}$
  and $\hat{I}_{\mathrm{phys}}$ is the identity on $\Hilb_{\mathrm{phys}}$.
\end{enumerate}
\end{definition}

\begin{theorem}[Commutation of Categorical and Physical Observables]
\label{thm:commutation}
Categorical observables commute with physical observables:
\begin{equation}
[\Ocat, \Ophys] = 0.
\label{eq:commutation}
\end{equation}
\end{theorem}

\begin{proof}
By the tensor product structure:
\begin{align}
\Ocat \, \Ophys
&= (\hat{I}_{\mathrm{phys}} \otimes \hat{B}_{\mathrm{cat}})
   (\hat{A}_{\mathrm{phys}} \otimes \hat{I}_{\mathrm{cat}})
\nonumber\\
&= (\hat{I}_{\mathrm{phys}} \hat{A}_{\mathrm{phys}})
   \otimes (\hat{B}_{\mathrm{cat}} \hat{I}_{\mathrm{cat}})
\nonumber\\
&= \hat{A}_{\mathrm{phys}}
   \otimes \hat{B}_{\mathrm{cat}}.
\label{eq:comm-forward}
\end{align}
Similarly:
\begin{align}
\Ophys \, \Ocat
&= (\hat{A}_{\mathrm{phys}} \otimes \hat{I}_{\mathrm{cat}})
   (\hat{I}_{\mathrm{phys}} \otimes \hat{B}_{\mathrm{cat}})
\nonumber\\
&= \hat{A}_{\mathrm{phys}}
   \otimes \hat{B}_{\mathrm{cat}}.
\label{eq:comm-reverse}
\end{align}
Therefore $\Ocat \, \Ophys = \Ophys \, \Ocat$, i.e.,
$[\Ocat, \Ophys] = 0$. $\qedhere$
\end{proof}

\begin{corollary}[Zero Backaction]
\label{cor:zero-backaction}
Categorical observation does not disturb the physical state. The
physical state before and after a categorical measurement is
identical. This is in stark contrast to quantum measurement on the
physical factor $\Hilb_{\mathrm{phys}}$, which generically causes
state collapse \citep{vonneumann1932}. The categorical observation
acts on a separate tensor factor and therefore commutes with all
physical dynamics.
\end{corollary}

\begin{remark}
The zero-backaction property is essential for the observation
apparatus interpretation. An instrument that disturbs what it
measures would require corrections; a categorical observer does not.
The fragment shader reads state and writes observation, leaving the
state unchanged. This is guaranteed by the tensor product structure.
\end{remark}


% ═════════════════════════════════════════════════════════════════════
%              PART II: THE RENDERING–MEASUREMENT IDENTITY
% ═════════════════════════════════════════════════════════════════════

\section{The Observation Operation}
\label{sec:observation}

We now formalize what it means to ``observe'' a categorical state.

\begin{definition}[Observation Operation]
\label{def:observation}
An \emph{observation} is a map
\begin{equation}
\Obs : \mathcal{S} \to \Tex,
\label{eq:observation-map}
\end{equation}
where $\mathcal{S} = [0,1]^3$ is the S-entropy coordinate space
(Definition~\ref{def:sentropy}) and
$\Tex = \R^{W \times H \times 4}$ is the space of RGBA textures
of resolution $W \times H$.

The observation $\Obs$ maps the categorical state --- compactly
represented as an S-entropy triple $(\Sk, \St, \Se)$ --- to a
texture $T \in \Tex$ such that each texel at coordinates $(u,v)$
with $u \in \{1,\ldots,W\}$, $v \in \{1,\ldots,H\}$ probes one
partition cell. The RGBA value $T(u,v) \in [0,1]^4$ at that texel
\emph{is} the observed categorical state of the corresponding cell.
\end{definition}

\begin{definition}[Partition-Texel Correspondence]
\label{def:partition-texel}
Define the correspondence map
$\pi : \{1,\ldots,W\} \times \{1,\ldots,H\} \to
\{C_\alpha\}_{\alpha=1}^{K}$
assigning each texel to a partition cell. For a $W \times H$ texture
observing $K \leq W \cdot H$ cells:
\begin{equation}
\pi(u,v) = C_\alpha, \quad \alpha = (v-1) \cdot W + u.
\label{eq:pi}
\end{equation}
The RGBA value at texel $(u,v)$ encodes the categorical state of
cell $C_{\pi(u,v)}$:
\begin{itemize}[leftmargin=*]
\item R channel: amplitude/intensity of the observed quantity
\item G channel: phase information (if applicable)
\item B channel: secondary observable (e.g., gradient magnitude)
\item A channel: confidence/weight
\end{itemize}
\end{definition}

\begin{proposition}[Observation Completeness]
\label{prop:completeness}
If $W \cdot H \geq K$, the observation $\Obs$ is \emph{complete}:
every partition cell is probed by at least one texel. The texture
$T = \Obs(\Sk, \St, \Se)$ records the full categorical state.
\end{proposition}

\begin{proof}
The map $\pi$ assigns $W \cdot H$ texels to $K$ cells. If
$W \cdot H \geq K$, every cell receives at least one texel (with
possible redundancy). The full categorical state is the assignment
of a value to each cell, which is exactly what the texture records.
\end{proof}


% ═════════════════════════════════════════════════════════════════════
\section{The Rendering--Measurement Identity}
\label{sec:rmi}

We now state and prove the main theorem of this paper.

\begin{definition}[Fragment Shader as Function Evaluator]
\label{def:fragment-shader}
A \emph{fragment shader} $F$ is a function
$F : [0,1]^2 \times \R^p \to [0,1]^4$
that takes normalized texel coordinates $(u',v') \in [0,1]^2$ and a
parameter vector $\bm{\theta} \in \R^p$ (uploaded as GPU uniforms)
and returns an RGBA color value. The GPU \emph{rendering operation}
$\Rend$ executes $F$ at every texel in parallel:
\begin{equation}
T = \Rend(F, \bm{\theta})
  \quad\Leftrightarrow\quad
  T(u,v) = F\!\left(\frac{u}{W}, \frac{v}{H},\, \bm{\theta}\right)
  \;\forall\,(u,v).
\label{eq:rendering}
\end{equation}
The GPU hardware guarantees that all texel evaluations are independent
and their order does not affect the result \citep{khronosglsl}.
\end{definition}

\begin{definition}[Categorical Measurement]
\label{def:categorical-measurement}
A \emph{categorical measurement} $\Meas$ on a system with S-entropy
state $\bm{s} = (\Sk, \St, \Se)$ is a map
$\Meas : \Obs \times \mathcal{S} \to \Tex$
that applies the observation operation $\Obs$ to the state $\bm{s}$
to produce a complete record of observed categorical values:
\begin{equation}
T = \Meas(\Obs, \bm{s})
  \quad\Leftrightarrow\quad
  T(u,v) = \Obs_{\pi(u,v)}(\bm{s})
  \;\forall\,(u,v),
\label{eq:measurement}
\end{equation}
where $\Obs_\alpha(\bm{s})$ is the observed value at partition cell
$C_\alpha$.
\end{definition}

\begin{theorem}[Rendering--Measurement Identity]
\label{thm:rmi}
For a fragment shader $F$ that implements the observation operation
$\Obs$ with parameter vector $\bm{\theta} = (\Sk, \St, \Se)$, the
rendering operation is identical to categorical measurement:
\begin{equation}
\boxed{
\Rend(F, \bm{s}) \equiv \Meas(\Obs, \bm{s}).
}
\label{eq:rmi}
\end{equation}
The texture $T$ output by the fragment shader is not a representation
of the measurement result --- it \emph{is} the measurement result.
\end{theorem}

\begin{proof}
We establish the identity by showing that the two operations are the
same map applied to the same input, producing the same output.

\textbf{Step 1: Input identity.} The rendering operation receives
the S-entropy coordinates $\bm{s} = (\Sk, \St, \Se)$ as GPU
uniform variables. The measurement operation receives the same
$\bm{s}$ as the state to be observed. Both start from the same
input.

\textbf{Step 2: Per-texel operation identity.} At each texel $(u,v)$:
\begin{itemize}[leftmargin=*]
\item The rendering operation evaluates
  $F(u/W,\, v/H,\, \bm{s})$, which computes the observation function
  at the partition cell corresponding to $(u,v)$.
\item The measurement operation evaluates
  $\Obs_{\pi(u,v)}(\bm{s})$, which is the observed value of the
  partition cell $C_{\pi(u,v)}$ given state $\bm{s}$.
\end{itemize}
Since $F$ implements $\Obs$ (by hypothesis), and the texel-to-cell
correspondence $\pi$ is encoded in the coordinate mapping $(u/W, v/H)$,
we have $F(u/W, v/H, \bm{s}) = \Obs_{\pi(u,v)}(\bm{s})$ for all
$(u,v)$.

\textbf{Step 3: Parallelism identity.} The rendering operation
evaluates all texels independently and in parallel. The measurement
operation probes all partition cells independently. By the commutativity
result (Theorem~\ref{thm:commutation}), the order of observation does
not affect the result. Both operations produce the same output
regardless of evaluation order.

\textbf{Step 4: Output identity.} Both operations write their
results to the same data structure: a $W \times H \times 4$ array
of floating-point values. The rendering operation writes to a GPU
framebuffer; the measurement operation records to a measurement
array. These are the same object.

\textbf{Step 5: Completeness.} By Proposition~\ref{prop:completeness},
if $W \cdot H \geq K$, both operations probe all partition cells.
The output is a complete record of all categorical observations.

\textbf{Conclusion.} The rendering operation and the measurement
operation have the same input ($\bm{s}$), the same per-element
function ($\Obs$ implemented by $F$), the same evaluation structure
(independent parallel evaluation), and the same output format
($W \times H \times 4$ texture). They are the same operation.
$\Rend(F, \bm{s}) \equiv \Meas(\Obs, \bm{s})$. \qedhere
\end{proof}

\begin{figure*}[!htbp]
    \centering
    \includegraphics[width=\textwidth]{figures/panel1_rendering_measurement_identity.png}
    \caption{\textbf{Rendering--Measurement Identity: Texture--State Isomorphism Verification.}
    (\textbf{A})~Three-dimensional scatter plot of categorical distance $d_{\mathrm{cat}}$, L2 texture distance $d_{L_2}$, and Dice coefficient for 200 random image pairs drawn from 50 synthetic BBBC039-like nuclei images. Points are coloured by Dice (viridis). The positive trend from lower-left to upper-right confirms that categorical distance (computed in partition coordinate space) and texture distance (computed in pixel space) are positively correlated---validating the Texture--State Isomorphism (Theorem~\ref{thm:isomorphism}).
    (\textbf{B})~Scatter plot of $d_{\mathrm{cat}}$ vs.\ $d_{L_2}$ with linear fit (red line). Pearson $r = 0.374$ ($p = 4.76 \times 10^{-8}$), demonstrating a statistically significant positive relationship between the two distance metrics. The moderate $r$ value (rather than $r \approx 1$) reflects the non-trivial, many-to-one nature of the partition encoding: multiple distinct pixel configurations can map to the same partition signature, introducing scatter around the linear trend. The key prediction---that $d_{\mathrm{cat}}$ and $d_{L_2}$ are monotonically related---is confirmed.
    (\textbf{C})~Histogram of categorical distances $d_{\mathrm{cat}}$ across all 200 pairs. The distribution is right-skewed with mean $\bar{d}_{\mathrm{cat}} = 0.051$ (red dashed), reflecting that most randomly selected image pairs have similar mean partition depths (both dominated by background pixels at low $n$). The tail extending to $d_{\mathrm{cat}} \approx 0.15$ corresponds to pairs with one high-nucleus-density image and one low-density image.
    (\textbf{D})~Histogram of L2 texture distances $d_{L_2}$ across all pairs. Mean $\bar{d}_{L_2} = 0.220$ (red dashed). The broader, more symmetric distribution compared to $d_{\mathrm{cat}}$ reflects the higher dimensionality of the pixel-space representation. The proportionality $d_{L_2} \approx 4.3 \cdot d_{\mathrm{cat}}$ (from the linear fit slope) provides the empirical scaling constant for the Texture--State Isomorphism on this dataset.}
    \label{fig:panel_identity}
\end{figure*}

% ── Figure 5 proposal ──────────────────────────────────────────────
% FIGURE 5: Two-panel diagram. Left: "Rendering" showing shader
% evaluating at each texel → texture. Right: "Measurement" showing
% observation at each partition cell → record. Central "≡" sign
% connecting them. Bottom: "Same operation, same result."

\begin{corollary}[Observation = Computation = Processing]
\label{cor:obs-comp}
Since categorical measurement \emph{is} computation (by the Triple
Equivalence, Theorem~\ref{thm:triple}, the categorical state encodes
the same information as the oscillatory and partitional states), and
rendering \emph{is} measurement (by Theorem~\ref{thm:rmi}), rendering
\emph{is} computation:
\begin{equation}
\text{Rendering} \equiv \text{Measurement}
  \equiv \text{Computation}.
\label{eq:triple-identity}
\end{equation}
The fragment shader does not ``accelerate'' computation.  It does not
``assist'' computation.  It \emph{is} computation.  There is no
separate computational step that the shader then visualizes.  The
single act of rendering performs the computation, produces the
measurement, and records the result simultaneously.
\end{corollary}

\begin{remark}[Paradigm Inversion]
Traditional pipeline: compute on CPU $\to$ transfer results $\to$
visualize on GPU. Observation pipeline: observe on GPU (= compute
= measure). There is no transfer because there is no separation.
The act of writing the pixel \emph{is} the act of computing the
answer.
\end{remark}


% ═════════════════════════════════════════════════════════════════════
\section{The Texture--State Isomorphism}
\label{sec:texture-state}

\begin{theorem}[Texture--State Isomorphism]
\label{thm:texture-state}
The texture $T = \Obs(\bm{s})$ output by the observation shader is
isomorphic to the categorical state $\sigma$ of the system. The
isomorphism $\Phi : \sigma \mapsto T$ preserves:
\begin{enumerate}[label=(\roman*),leftmargin=*]
\item \textbf{Metric structure:}
  $\dcat(\sigma_1, \sigma_2)
  = \lambda \cdot \|T_1 - T_2\|_{L^2}$
  for a normalization constant $\lambda > 0$.
\item \textbf{Entropy:} $S(\sigma) = S(T)$.
\item \textbf{Information:} $I(\sigma) = I(T)$.
\end{enumerate}
\end{theorem}

\begin{proof}
\textbf{(i) Metric preservation.}
The categorical distance (Definition~\ref{def:cat-distance}) is:
$\dcat(\sigma_1, \sigma_2) = \sum_\alpha |v_{1,\alpha} - v_{2,\alpha}|$,
where $v_{i,\alpha}$ is the categorical value of state $\sigma_i$ at
cell $C_\alpha$. Under the partition-texel correspondence $\pi$, this
maps to:
\begin{equation}
\dcat(\sigma_1, \sigma_2)
= \sum_{u,v} |T_1(u,v) - T_2(u,v)|
= \|T_1 - T_2\|_{L^1}.
\end{equation}
By equivalence of norms on finite-dimensional spaces, $\|T_1 - T_2\|_{L^1}
\leq \sqrt{W H} \cdot \|T_1 - T_2\|_{L^2} \leq W H \cdot
\|T_1 - T_2\|_{L^\infty}$. Since the texels encode the same values as
the categorical labels, the $L^1$ texture distance equals the
categorical distance exactly. The $L^2$ distance is related by the
Cauchy--Schwarz factor $\lambda = 1/\sqrt{WH}$.

\textbf{(ii) Entropy preservation.}
The Shannon entropy of the categorical state is:
$S(\sigma) = -\sum_\alpha p_\alpha \ln p_\alpha$.
The Shannon entropy of the texture, treating pixel values as outcomes
of a random variable over the partition, is identical because the
texture records exactly the same values:
$S(T) = -\sum_{u,v} p_{T(u,v)} \ln p_{T(u,v)} = S(\sigma)$.

\textbf{(iii) Information preservation.}
The mutual information between the state and any external variable
$Y$ is:
$I(\sigma; Y) = S(\sigma) - S(\sigma|Y) = S(T) - S(T|Y) = I(T; Y)$,
since $\sigma$ and $T$ carry identical content by the bijection
$\Phi$. \qedhere
\end{proof}

\begin{remark}
The Texture--State Isomorphism means that operating on the texture
\emph{is} operating on the categorical state. Computing an $L^2$
distance between two textures \emph{is} computing the categorical
distance between two states. This is not a proxy or approximation ---
it is exact.
\end{remark}


% ═════════════════════════════════════════════════════════════════════
%              PART III: O(1) MEMORY COMPLEXITY
% ═════════════════════════════════════════════════════════════════════

\section{The Storage Redundancy Theorem}
\label{sec:storage-redundancy}

\begin{theorem}[Storage Redundancy]
\label{thm:storage-redundancy}
If the observation apparatus (fragment shader $F$) resides in GPU
memory and the raw state description (S-entropy coordinates
$\bm{s} \in [0,1]^3$) is available, then storing the observation
result (texture $T$) is redundant: $T$ can be regenerated in
$\BigO(1)$ time (one GPU draw call).
\end{theorem}

\begin{proof}
By the Rendering--Measurement Identity (Theorem~\ref{thm:rmi}):
$T = \Rend(F, \bm{s})$.

This relation is deterministic: the same $F$ and $\bm{s}$ always
produce the same $T$ (the GPU rendering pipeline is
deterministic for fixed shader code and inputs
\citep{khronosglsl,khronosvulkan}).

The time to regenerate $T$ is one draw call: upload $\bm{s}$ to a
GPU uniform (constant time, 12 bytes for three floats), dispatch
the fragment shader (constant time, one draw call over a fullscreen
quad), and read the result from the framebuffer (constant time for
fixed texture resolution). None of these steps depends on any
external parameter such as database size $N$.

Therefore, storing $T$ is unnecessary. It can be regenerated on
demand in $\BigO(1)$ time. Any system that stores $N$ observation
results uses $\BigO(N)$ memory redundantly, since all $N$ can be
regenerated from $N$ S-entropy triples (12 bytes each) plus the
single apparatus $F$ (already in GPU memory). \qedhere
\end{proof}

\begin{corollary}[Memory Reduction]
\label{cor:memory-reduction}
The standard approach stores $N$ observation results:
$M_{\mathrm{standard}} = N \cdot W \cdot H \cdot 4 \cdot
\text{sizeof(float)} = \BigO(N)$.
The observation approach stores one apparatus plus streams states:
$M_{\mathrm{obs}} = M_{\mathrm{shader}} + c = \BigO(1)$,
where $c$ is a constant (the working set of textures).
The memory reduction ratio is $N / 1 = N$.
\end{corollary}


% ═════════════════════════════════════════════════════════════════════
\section{The Streaming Observation Protocol}
\label{sec:streaming}

\begin{definition}[Streaming Observation Protocol]
\label{def:streaming-protocol}
The streaming observation protocol processes a database of $N$ items
with $\BigO(1)$ GPU memory:
\end{definition}

\begin{algorithm}[H]
\caption{Streaming Observation Protocol}
\label{alg:streaming}
\begin{algorithmic}[1]
\Require Database $\mathcal{D} = \{\bm{s}_1, \ldots, \bm{s}_N\}$
  of S-entropy triples
\Require Query S-entropy triple $\bm{s}_q$
\Require Observation apparatus $F$ (shader programs)
\Require Threshold $\tau$
\Ensure Best match index $i^*$ and score $d^*$
\State \textbf{Phase 0: Load apparatus.}
  Compile and link shader programs $F$ on GPU.
  \Comment{Once, $\sim$10 MB}
\State Render query observation:
  $T_q \gets \Rend(F, \bm{s}_q)$
  \Comment{$\BigO(1)$, one draw call}
\State $d^* \gets \infty$; $i^* \gets \textsc{null}$
\For{$i = 1$ to $N$}
  \Comment{Stream items one by one}
  \State Upload $\bm{s}_i$ to GPU uniform
    \Comment{12 bytes}
  \State Render item observation:
    $T_i \gets \Rend(F, \bm{s}_i)$
    \Comment{$\BigO(1)$}
  \State Compute interference:
    $T_{\mathrm{int}} \gets \Rend(F_{\mathrm{int}}, T_q, T_i)$
    \Comment{$\BigO(1)$}
  \State Extract score:
    $d_i \gets \text{ReadBack}(T_{\mathrm{int}})$
    \Comment{$\BigO(1)$}
  \If{$d_i < d^*$}
    \State $d^* \gets d_i$; $i^* \gets i$
  \EndIf
  \State \textbf{Discard} $T_i$ and $T_{\mathrm{int}}$
    \Comment{Memory freed}
\EndFor
\State \Return $(i^*, d^*)$
\end{algorithmic}
\end{algorithm}

\begin{proposition}[GPU Memory Invariance]
\label{prop:memory-invariance}
During the execution of Algorithm~\ref{alg:streaming}, the GPU
memory usage is constant:
\begin{equation}
\Mgpu = \Mshader + 3 \cdot W \cdot H \cdot 4 \cdot
  \text{sizeof(float)},
\label{eq:Mgpu}
\end{equation}
where the factor $3$ accounts for three textures: query ($T_q$),
current item ($T_i$), and interference buffer ($T_{\mathrm{int}}$).
This is independent of the database size $N$.
\end{proposition}

\begin{proof}
At any point during the loop, the GPU holds: (i)~the shader programs
$F$ and $F_{\mathrm{int}}$ ($\Mshader$), (ii)~the query texture $T_q$
(retained for all comparisons), (iii)~the current item texture $T_i$
(overwritten each iteration), and (iv)~the interference texture
$T_{\mathrm{int}}$ (overwritten each iteration). No other GPU memory
is allocated. The total is $\Mshader + 3 \cdot W \cdot H \cdot 4
\cdot \text{sizeof(float)}$. Since $W$, $H$, and $\Mshader$ are
constants independent of $N$, the GPU memory is $\BigO(1)$.
\end{proof}


% ═════════════════════════════════════════════════════════════════════
\section{Memory Bound Theorem}
\label{sec:memory-bound}

\begin{theorem}[Explicit Memory Bound]
\label{thm:memory-bound}
For texture resolution $W = H = 256$ and 32-bit floating point:
\begin{align}
\Mgpu &= \Mshader + 3 \times 256 \times 256 \times 16
  \;\text{bytes} \nonumber\\
&\approx 10\,\text{MB} + 3 \times 1{,}048{,}576
  \;\text{bytes} \nonumber\\
&\approx 10\,\text{MB} + 3.15\,\text{MB} \nonumber\\
&\approx 13.15\,\text{MB}.
\label{eq:explicit-bound}
\end{align}
\end{theorem}

\begin{proof}
Shader programs (vertex shader, observation shader, interference
shader, quality metric shader, display shader) together with their
compiled binaries, uniform buffers, and ancillary state occupy
approximately 10\,MB on typical GPU drivers
\citep{khronosglsl,khronosvulkan}. Each texture of resolution
$256 \times 256$ with 4 channels (RGBA) of 32-bit floats occupies
$256 \times 256 \times 4 \times 4 = 1{,}048{,}576$ bytes
$\approx 1.05$\,MB. Three such textures total $\approx 3.15$\,MB.
The sum is $\approx 13.15$\,MB. \qedhere
\end{proof}

\begin{remark}[Comparison with Standard Approaches]
Consider a vector database storing $N = 10^8$ items with feature
dimension $d = 768$ (a typical BERT embedding dimension) in 32-bit
float:
\begin{equation}
M_{\mathrm{VDB}} = N \cdot d \cdot 4 = 10^8 \times 768 \times 4
  \approx 307\,\text{GB}.
\end{equation}
The observation approach requires $\sim$13\,MB regardless of $N$.
The ratio is $307\,\text{GB} / 13\,\text{MB} \approx 23{,}600\times$.
\end{remark}

\begin{figure*}[!htbp]
    \centering
    \includegraphics[width=\textwidth]{figures/panel2_memory_scaling.png}
    \caption{\textbf{$\mathcal{O}(1)$ Memory Scaling: Observation vs.\ Standard Approaches.}
    (\textbf{A})~Three-dimensional surface plot of GPU memory usage as a function of database size $N$ and approach type. The observation approach (blue) forms a flat plane at $\sim$13\,MB regardless of $N$. The standard vector-database approach (orange) forms a steeply rising surface growing as $\mathcal{O}(N \cdot d)$ with $d = 128$. The surfaces intersect at $N \approx 250$, beyond which the observation approach is strictly more memory-efficient.
    (\textbf{B})~Log-log plot of GPU memory vs.\ database size $N$ for $N \in \{10^2, 10^3, 10^4, 10^5, 10^6\}$. The observation approach (blue circles, solid) produces a horizontal line at $13.15$\,MB---confirming the Memory Bound Theorem (Theorem~\ref{thm:memory_bound}). The standard approach (orange squares, dashed) grows linearly with slope~1 on the log-log axes, reaching $48.8$\,MB at $N = 10^5$. The constant-vs-linear divergence is the central prediction of the Storage Redundancy Theorem.
    (\textbf{C})~Bar chart comparing memory at $N = 10^5$: observation approach requires $13.15$\,MB; standard approach requires $48.8$\,MB---a $3.7\times$ advantage. At $N = 10^8$ (not shown), the ratio grows to $\sim$$23{,}000\times$, as the standard approach requires $\sim$307\,GB while the observation approach remains at $13.15$\,MB.
    (\textbf{D})~Bar chart of memory ratio (standard / observation) as a function of database size. The ratio grows linearly with $N$, confirming the $\mathcal{O}(N)$ vs.\ $\mathcal{O}(1)$ scaling difference. At $N = 10^6$, the observation approach uses $37\times$ less GPU memory. This linear growth of the advantage with database size is the practical consequence of the Storage Redundancy Theorem: larger databases benefit more from the observation paradigm.}
    \label{fig:panel_memory}
\end{figure*}


% ═════════════════════════════════════════════════════════════════════
\section{Batch Observation for Throughput}
\label{sec:batch}

\begin{definition}[Batched Streaming]
\label{def:batch}
In batched streaming, $B$ items are observed in parallel rather than
sequentially. The GPU allocates $B$ item textures and processes them
simultaneously, then discards the batch and proceeds to the next.
\end{definition}

\begin{proposition}[Batch Memory and Throughput]
\label{prop:batch}
For batch size $B$:
\begin{align}
M_{\mathrm{batch}} &= \Mshader + (B + 2) \cdot W \cdot H
  \cdot 4 \cdot \text{sizeof(float)}, \\
\text{Throughput} &= \frac{B}{T_{\mathrm{render}}},
\label{eq:batch}
\end{align}
where $T_{\mathrm{render}}$ is the render time for one batch.
\end{proposition}

\begin{example}[Practical Batch on 4\,GB Integrated GPU]
For $B = 2000$, $W = H = 256$:
\begin{align}
M_{\mathrm{batch}} &= 10\,\text{MB}
  + 2002 \times 256 \times 256 \times 16\,\text{bytes}
\nonumber\\
&\approx 10\,\text{MB} + 2002 \times 1.05\,\text{MB}
\nonumber\\
&\approx 2.1\,\text{GB}.
\end{align}
This fits within the 4\,GB memory of a typical integrated GPU.
At a render time of $\sim$50\,$\mu$s per batch (the fragment shader
evaluates all $2000 \times 256 \times 256 \approx 1.3 \times 10^8$
texels in parallel), the throughput is:
\begin{equation}
\text{Throughput} \approx \frac{2000}{50\,\mu\text{s}}
  = 40 \times 10^6 \;\text{observations/second}.
\end{equation}
\end{example}


% ═════════════════════════════════════════════════════════════════════
%        PART IV: GPU-SUPERVISED TRAINING (NO HUMAN LABELS)
% ═════════════════════════════════════════════════════════════════════

\section{Physical Observable Extraction}
\label{sec:observables}

\begin{theorem}[Physical Observable Extraction]
\label{thm:observables}
The observation texture $T$ produced by the fragment shader contains
objective, deterministic physical quality metrics that can be
extracted without human judgment. The following five observables are
measurable from any observation texture:
\end{theorem}

\begin{definition}[GPU-Measurable Physical Observables]
\label{def:gpu-observables}
Given an observation texture $T \in \R^{W \times H \times 4}$:

\textbf{(a) Partition Sharpness.}
\begin{equation}
P_{\mathrm{sharp}} = \frac{1}{N_{\mathrm{texels}}}
  \sum_{u,v} |\nabla T(u,v)|,
\label{eq:sharpness}
\end{equation}
where $\nabla T(u,v)$ is the discrete gradient (Sobel or central
differences) at texel $(u,v)$. Higher values indicate better-resolved
partition boundaries. This is measurable because it depends only on
finite differences of pixel values.

\textbf{(b) Observation Noise.}
\begin{equation}
N_{\mathrm{obs}} = \frac{\sum_{|\bm{k}| > k_c}
  |\hat{T}(\bm{k})|^2}{\sum_{\bm{k}} |\hat{T}(\bm{k})|^2},
\label{eq:noise}
\end{equation}
where $\hat{T}$ is the discrete Fourier transform of $T$ and $k_c$
is a cutoff frequency. Lower values indicate cleaner observation.
This is computable via FFT on the GPU texture.

\textbf{(c) Phase Coherence.}
\begin{equation}
C_{\mathrm{phase}} = \frac{1}{N_{\mathrm{pairs}}}
  \left|\sum_{\langle u,v \rangle \sim \langle u',v' \rangle}
  \exp\!\big(i\,\phi(u,v) - i\,\phi(u',v')\big)\right|,
\label{eq:coherence}
\end{equation}
where $\phi(u,v)$ is the phase angle extracted from the texture value
at $(u,v)$ and the sum runs over neighboring texel pairs. Higher
values indicate more coherent structure across partition cells.

\textbf{(d) Interference Visibility.}
\begin{equation}
V_{\mathrm{int}} = \frac{I_{\max} - I_{\min}}{I_{\max} + I_{\min}},
\label{eq:visibility}
\end{equation}
where $I_{\max}$ and $I_{\min}$ are the maximum and minimum
intensities of the interference pattern between the observation
texture and a reference. This is the Michelson contrast
\citep{michelson1891,bornwolf1999}. Higher values indicate
stronger distinguishable signal.

\textbf{(e) Multi-Resolution Consistency.}
\begin{equation}
R_{\mathrm{multi}} = 1 - \frac{|T_\sigma - T_{2\sigma}|}{|T_\sigma|},
\label{eq:multires}
\end{equation}
where $T_\sigma$ is the texture blurred with a Gaussian of standard
deviation $\sigma$ and $T_{2\sigma}$ is blurred with $2\sigma$.
Higher values indicate stability across spatial scales: the
observation does not change qualitatively when viewed at different
resolutions.
\end{definition}

\begin{proof}[Proof of Theorem~\ref{thm:observables}]
Each observable is a deterministic function of the texture $T$:
\begin{itemize}[leftmargin=*]
\item $P_{\mathrm{sharp}}$ depends only on finite differences of
  pixel values --- no judgment required.
\item $N_{\mathrm{obs}}$ depends only on the Fourier transform ---
  a linear operator on the pixel array.
\item $C_{\mathrm{phase}}$ depends only on pixel values at
  neighboring texels --- a local computation.
\item $V_{\mathrm{int}}$ depends only on the maximum and minimum
  intensity of a derived texture --- min/max operations.
\item $R_{\mathrm{multi}}$ depends only on Gaussian convolution
  of the texture at two scales --- a standard linear filter.
\end{itemize}
All five are computable in a single additional render pass (or at most
two passes for the FFT). None requires human annotation, subjective
judgment, or external data. They are physical observables of the
observation texture, analogous to measuring the visibility of an
interference fringe in an optics experiment
\citep{bornwolf1999,michelson1891}. \qedhere
\end{proof}

% ── Figure 6 proposal ──────────────────────────────────────────────
% FIGURE 6: Five panels showing each observable extracted from an
% example observation texture: (a) gradient map for sharpness,
% (b) high-frequency residual for noise, (c) phase coherence map,
% (d) interference fringe pattern with visibility annotation,
% (e) multi-resolution comparison.

\begin{figure*}[!htbp]
    \centering
    \includegraphics[width=\textwidth]{figures/panel3_physical_observables.png}
    \caption{\textbf{Physical Observable Extraction and Quality Correlation.}
    (\textbf{A})~Three-dimensional scatter plot of the three primary GPU-measured physical observables---partition sharpness $P_{\mathrm{sharp}}$, observation noise $N_{\mathrm{obs}}$, and phase coherence $C_{\mathrm{phase}}$---for 50 synthetic images, coloured by segmentation Dice coefficient (proxy for observation quality). Higher-Dice images (green--yellow) cluster in the high-sharpness, low-noise, high-coherence region, confirming that the physical observables are informative predictors of observation quality without requiring ground-truth labels.
    (\textbf{B})~Scatter plot of partition sharpness $P_{\mathrm{sharp}}$ vs.\ Dice coefficient with linear trend line. Higher sharpness corresponds to better-resolved partition boundaries, which in turn enables more accurate segmentation. The positive correlation validates sharpness as an objective, label-free quality metric: images with sharp partition boundaries are images where the observation apparatus (fragment shader) has successfully resolved the categorical structure.
    (\textbf{C})~Scatter plot of observation noise $N_{\mathrm{obs}}$ vs.\ Dice. Lower noise corresponds to higher Dice, as expected: noisy observations obscure partition boundaries. The negative trend confirms that $N_{\mathrm{obs}}$ provides a complementary quality signal to $P_{\mathrm{sharp}}$---the two together capture both the strength of the signal (sharpness) and the magnitude of the interference (noise).
    (\textbf{D})~Scatter plot of phase coherence $C_{\mathrm{phase}}$ vs.\ Dice. Phase coherence measures the consistency of the oscillatory phase field across neighbouring partition cells. Higher coherence indicates spatially smooth structure (biological tissue), while low coherence indicates fragmented or noisy observations. The positive trend confirms $C_{\mathrm{phase}}$ as the third independent component of the GPU-supervised training signal (Eq.~\ref{eq:loss}).}
    \label{fig:panel_observables}
\end{figure*}

% ═════════════════════════════════════════════════════════════════════
\section{GPU-Supervised Training Framework}
\label{sec:training}

\begin{definition}[Training Architecture]
\label{def:training-arch}
The GPU-supervised training architecture comprises four components
with distinct roles:
\end{definition}

\subsection{Component 1: Compiled Probe (Neural, Trainable)}

The compiled probe is a neural network with parameters $\theta$:
\begin{equation}
\text{Probe}_\theta : \mathcal{X} \to \mathcal{A},
\label{eq:probe}
\end{equation}
where $\mathcal{X}$ is the input space (natural language query,
domain data, or raw signal) and $\mathcal{A}$ is the space of
\emph{symbolic operation sequences} --- structured instructions
specifying which observations to perform.

\textbf{Architecture:} A frozen pre-trained language model base
(e.g., a 7B parameter model) with LoRA (Low-Rank Adaptation
\citep{hu2022}) adapters. Only the LoRA adapters ($\sim$10M
parameters) are trainable; the base model is frozen.

\textbf{Key property:} This is the \emph{only} component with
learnable parameters. All other components are deterministic.

\subsection{Component 2: Symbolic Executor (Deterministic)}

\begin{equation}
\text{Exec} : \mathcal{A} \to \mathcal{G},
\label{eq:executor}
\end{equation}
where $\mathcal{G}$ is the space of GPU dispatch commands (shader
selection, uniform values, render targets). This is a pure function
with no learnable parameters. It translates the operation sequence
from the probe into concrete GPU calls.

\subsection{Component 3: GPU Observation Apparatus (Physical)}

\begin{equation}
\text{GPU} : \mathcal{G} \to (\Tex, \bm{q}),
\label{eq:gpu-apparatus}
\end{equation}
where $\Tex$ is the observation texture and $\bm{q} = (P_{\mathrm{sharp}},
N_{\mathrm{obs}}, C_{\mathrm{phase}}, V_{\mathrm{int}},
R_{\mathrm{multi}})$ is the vector of physical observables. The GPU
executes the fragment shaders, producing both the observation and
its quality metrics. There are no learnable parameters. The GPU
acts as a \emph{black-box oracle}: given an input, it produces a
deterministic output.

\subsection{Component 4: Training Signal Extractor (Deterministic)}

\begin{equation}
\text{Loss} : \bm{q} \to \R,
\label{eq:loss-extractor}
\end{equation}
a deterministic function mapping the physical observables to a scalar
loss value. No learnable parameters.

% ── Figure 7 proposal ──────────────────────────────────────────────
% FIGURE 7: Training architecture diagram. Four boxes connected
% by arrows: Probe(θ) → Executor → GPU Oracle → Loss.
% Only the Probe box is shaded (trainable). Gradient arrow flows
% back from Loss to Probe.

\begin{figure*}[!htbp]
    \centering
    \includegraphics[width=\textwidth]{figures/panel4_gpu_supervised_training.png}
    \caption{\textbf{GPU-Supervised Training: Convergence Without Human Labels.}
    (\textbf{A})~Three-dimensional surface plot of the composite loss landscape as a function of two probe parameters $w_1$ and $w_2$. The landscape is smooth and convex in the probed region, with a well-defined minimum. This smoothness is a prerequisite for gradient-based optimization and validates the Black-Box Oracle Convergence Theorem (Theorem~\ref{thm:convergence}): even though the GPU observation apparatus is a non-differentiable black box, the loss landscape induced by the physical observables is smooth enough for effective optimization.
    (\textbf{B})~Training loss curves over 100 iterations for two approaches: GPU-supervised training (blue, using only physical observables $P_{\mathrm{sharp}}$, $N_{\mathrm{obs}}$, $C_{\mathrm{phase}}$ as loss terms) and supervised baseline (orange, using ground-truth Dice as loss). Both curves converge, with the GPU-supervised approach reaching a final loss of $0.050$ and the supervised baseline reaching $0.009$. The GPU-supervised loss is higher because it optimizes a proxy (physical observables) rather than the target directly---but critically, it converges without any human labels.
    (\textbf{C})~Bar chart comparing final loss values: GPU-supervised ($0.050$) vs.\ supervised ($0.009$). The $5.8\times$ gap represents the cost of label-free training. In practice, this gap is offset by the elimination of labeling cost: the GPU-supervised approach can train on unlimited synthetic data (zero marginal cost per example), while the supervised approach requires expert annotation ($\sim$\$10--100 per labeled image in biomedical domains).
    (\textbf{D})~Observable improvement over training iterations: partition sharpness and phase coherence increase monotonically as the probe learns to generate better observation sequences, while noise decreases. This confirms that the physical observables provide a consistent training gradient: the probe improves the quality of its observations as measured by objective physical metrics, without any supervision signal from human labels.}
    \label{fig:panel_training}
\end{figure*}

% ═════════════════════════════════════════════════════════════════════
\section{The Composite Loss Function}
\label{sec:loss}

\begin{definition}[Composite Loss]
\label{def:composite-loss}
The training loss is:
\begin{align}
L(\theta) &= \lambda_1 \cdot L_{\mathrm{correct}}
  + \lambda_2 \cdot (1 - P_{\mathrm{sharp}})
  + \lambda_3 \cdot N_{\mathrm{obs}} \nonumber\\
&\quad + \lambda_4 \cdot (1 - C_{\mathrm{phase}})
  + \lambda_5 \cdot (1 - V_{\mathrm{int}}) \nonumber\\
&\quad + \lambda_6 \cdot (1 - R_{\mathrm{multi}})
  + \lambda_7 \cdot \frac{|\text{ops}|}{\text{max\_ops}},
\label{eq:loss}
\end{align}
where:
\begin{itemize}[leftmargin=*]
\item $L_{\mathrm{correct}}$: supervised loss when ground truth is
  available (cross-entropy, MSE, etc.); set to $0$ when no labels
  exist.
\item $P_{\mathrm{sharp}}, N_{\mathrm{obs}}, C_{\mathrm{phase}},
  V_{\mathrm{int}}, R_{\mathrm{multi}}$: GPU-measured physical
  observables (Definition~\ref{def:gpu-observables}).
\item $|\text{ops}|/\text{max\_ops}$: operation count regularizer,
  encouraging the probe to produce concise operation sequences.
\item $\lambda_1, \ldots, \lambda_7 \geq 0$: weighting
  coefficients.
\end{itemize}
\end{definition}

\begin{proposition}[Label-Free Training]
\label{prop:label-free}
Terms $2$--$7$ of the composite loss are always available, even
without labeled data. Training can proceed with $\lambda_1 = 0$
(no supervised signal). The GPU observables provide a
self-consistent training signal that drives the probe to produce
operation sequences yielding high-quality observations: sharp
partitions, low noise, coherent phase, high interference
visibility, and multi-resolution stability.
\end{proposition}

\begin{proof}
Each of terms $2$--$7$ is a function only of the observation texture
$T$ and the operation sequence length $|\text{ops}|$, both of which
are produced by the pipeline regardless of whether ground-truth labels
exist. The gradient $\nabla_\theta L$ can be computed as long as the
dependence of $L$ on $\theta$ is differentiable or estimable (see
\S\ref{sec:convergence}). No labeled data is required. \qedhere
\end{proof}

\begin{remark}
This is the key insight of GPU-supervised training: the physics of
the observation process --- the quality of the measurement --- provides
an objective, label-free training signal. A ``good'' observation has
sharp boundaries, low noise, coherent phases, high visibility, and
scale-stable structure. A ``bad'' observation has the opposite. The
GPU measures this directly, as an experimentalist would evaluate the
quality of an interference pattern.
\end{remark}


% ═════════════════════════════════════════════════════════════════════
\section{The Black-Box Oracle Convergence Theorem}
\label{sec:convergence}

\begin{theorem}[Black-Box Oracle Convergence]
\label{thm:convergence}
The training loop
\begin{equation}
\theta_{t+1} = \theta_t
  - \eta \cdot \hat{\nabla}_\theta L(\theta_t),
\label{eq:update}
\end{equation}
where $\hat{\nabla}_\theta L$ is an estimated gradient, converges to
a local minimum of $L$ under the following conditions:
\begin{enumerate}[label=(\roman*),leftmargin=*]
\item The probe $\mathrm{Probe}_\theta$ is a smooth parameterized
  function (neural network with smooth activations).
\item The executor is a deterministic pure function.
\item The GPU oracle is deterministic (same input $\to$ same output).
\item The loss $L$ is a smooth function of the observables.
\item The gradient estimator $\hat{\nabla}_\theta L$ is unbiased
  with bounded variance.
\end{enumerate}
\end{theorem}

\begin{proof}
The overall pipeline is:
\begin{equation}
L(\theta) = \text{Loss}\big(\text{GPU}\big(
  \text{Exec}\big(\text{Probe}_\theta(x)\big)\big)\big).
\label{eq:pipeline}
\end{equation}

\textbf{Case 1: Continuous probe output.} If the probe produces
continuous outputs (e.g., S-entropy coordinates directly), then
the composition is differentiable (smooth probe composed with
deterministic smooth executor, deterministic GPU, smooth loss). The
standard gradient $\nabla_\theta L$ exists and stochastic gradient
descent converges to a local minimum by the Robbins--Monro theorem,
given standard conditions on the learning rate
($\sum \eta_t = \infty$, $\sum \eta_t^2 < \infty$).

\textbf{Case 2: Discrete probe output.} If the probe produces
discrete operation sequences (typical for symbolic executors), the
pipeline is non-differentiable at the discrete step. We employ the
REINFORCE estimator \citep{williams1992}:
\begin{equation}
\hat{\nabla}_\theta L = \mathbb{E}_{a \sim \pi_\theta}
  \big[(L(a) - b) \cdot \nabla_\theta \log \pi_\theta(a|x)\big],
\label{eq:reinforce}
\end{equation}
where $\pi_\theta(a|x)$ is the probe's policy (probability of
generating action sequence $a$ given input $x$) and $b$ is a
baseline for variance reduction. The REINFORCE estimator is
unbiased \citep{williams1992,sutton1999}, satisfying condition (v).

Alternatively, the straight-through estimator
\citep{bengio2013estimating} treats discrete operations as continuous
in the backward pass, providing a biased but low-variance gradient
estimate that works well in practice.

In either case, the gradient estimator satisfies the conditions of
the stochastic approximation convergence theorem: unbiased (or
asymptotically unbiased) gradient estimates with diminishing learning
rate guarantee convergence to a local minimum.

\textbf{Determinism of the oracle.} The GPU oracle is deterministic
(condition (iii)), which means the loss landscape $L(\theta)$ is fixed
--- it does not shift between evaluations. This is a crucial
advantage over stochastic oracles: the same probe output always
produces the same loss, eliminating one source of variance. \qedhere
\end{proof}


% ═════════════════════════════════════════════════════════════════════
%            PART V: INTEGRATED GPU SUFFICIENCY
% ═════════════════════════════════════════════════════════════════════

\section{The Hardware Requirement Theorem}
\label{sec:hardware}

\begin{theorem}[Integrated GPU Sufficiency]
\label{thm:hardware}
The observation-based computation pipeline requires:
\begin{enumerate}[label=(\roman*),leftmargin=*]
\item \textbf{GPU memory:} $\sim$25\,MB working set.
\item \textbf{GPU compute:} $\sim$1--2\,TFLOPS.
\item \textbf{CPU memory:} $\sim$600\,MB (probe model + streaming
  buffer).
\end{enumerate}
These requirements are met by any modern integrated GPU. No discrete
GPU, cloud infrastructure, or datacenter is required.
\end{theorem}

\begin{proof}
We derive each requirement and compare to integrated GPU
specifications.

\textbf{(i) GPU memory.}
From Theorem~\ref{thm:memory-bound}: the streaming observation
protocol requires $\sim$13\,MB. Adding a margin for the quality
metric extraction shaders, intermediate buffers, and driver overhead:
\begin{align}
\Mgpu &= 13\,\text{MB (streaming)}
  + 3\,\text{MB (metric textures)} \nonumber\\
&\quad + 5\,\text{MB (quality extraction)}
  + 4\,\text{MB (driver overhead)} \nonumber\\
&= 25\,\text{MB}.
\end{align}

Integrated GPU specifications \citep{inteluhd770,applem2,amdradeon680m}:
\begin{center}
\begin{tabular}{lrr}
\toprule
\textbf{GPU} & \textbf{Shared Memory} & \textbf{Available} \\
\midrule
Intel UHD 770 & up to 64\,GB & $\geq$2\,GB \\
Apple M2 (10-core) & unified 8--24\,GB & $\geq$4\,GB \\
AMD Radeon 680M & shared 16--32\,GB & $\geq$2\,GB \\
\bottomrule
\end{tabular}
\end{center}
All exceed 25\,MB by a factor of $80$--$160\times$.

\textbf{(ii) GPU compute.}
Each observation requires evaluating the fragment shader at
$256 \times 256 = 65{,}536$ texels. Each texel evaluation involves
$\sim$50--200 floating-point operations (trigonometric functions,
multiplications, additions for the partition observation). Total per
observation:
\begin{equation}
\text{FLOPs/obs} = 65{,}536 \times 200
  \approx 1.3 \times 10^7 \;\text{FLOPs}.
\end{equation}
For $40 \times 10^6$ observations per second (batched):
\begin{equation}
\text{TFLOPS} = 40 \times 10^6 \times 1.3 \times 10^7
  / 10^{12} \approx 0.5\,\text{TFLOPS}.
\end{equation}

Integrated GPU capabilities:
\begin{center}
\begin{tabular}{lr}
\toprule
\textbf{GPU} & \textbf{TFLOPS (FP32)} \\
\midrule
Intel UHD 770 & $\sim$0.8 \\
Apple M2 (10-core GPU) & $\sim$3.6 \\
AMD Radeon 680M & $\sim$3.4 \\
\bottomrule
\end{tabular}
\end{center}
All exceed the requirement by $1.6$--$7\times$.

\textbf{(iii) CPU memory.}
The compiled probe ($\sim$10M LoRA parameters at 16-bit:
$10 \times 10^6 \times 2 = 20$\,MB) plus the frozen base model
quantized to 4-bit ($\sim$500\,MB for a 7B model) plus streaming
buffers ($\sim$80\,MB for a batch of S-entropy triples):
\begin{equation}
M_{\mathrm{CPU}} \approx 500 + 20 + 80 = 600\,\text{MB}.
\end{equation}
This is well within the 8--32\,GB RAM of any modern laptop.

\textbf{Comparison with standard ML.} Transformer inference for a
7B model requires 14--28\,GB GPU memory (FP16/FP32), plus
$\sim$100\,TFLOPS for real-time generation. Vector database retrieval
over $10^8$ items requires $\sim$300\,GB memory. The observation
approach requires $\sim$25\,MB GPU memory and $\sim$0.5\,TFLOPS.
The reduction is three to four orders of magnitude. \qedhere
\end{proof}

% ── Figure 8 proposal ──────────────────────────────────────────────
% FIGURE 8: Bar chart comparing hardware requirements:
% (a) GPU memory: observation (25 MB) vs. transformer (20 GB) vs.
%     vector DB (300 GB).
% (b) Compute: observation (0.5 TFLOPS) vs. transformer (100 TFLOPS).
% Log scale on y-axis. Labels on bars.


% ═════════════════════════════════════════════════════════════════════
\section{The Zero-Copy Advantage}
\label{sec:zerocopy}

\begin{theorem}[Integrated GPU Advantage]
\label{thm:zerocopy}
For the observation workload, integrated GPUs (which share main memory
with the CPU) have a structural advantage over discrete GPUs up to
a crossover database size $N_c$.
\end{theorem}

\begin{proof}
We analyze the data transfer costs.

\textbf{Per-observation data transfer.}
\begin{itemize}[leftmargin=*]
\item \textbf{Upload:} S-entropy coordinates --- 3 floats = 12 bytes.
\item \textbf{Readback:} Quality metrics --- 5 floats = 20 bytes,
  or a single scalar score = 4 bytes.
\end{itemize}
Total per observation: $\sim$16--32 bytes of transfer.

\textbf{Discrete GPU (PCIe transfer):}
PCIe 4.0 $\times$16 bandwidth: $\sim$32\,GB/s. Latency per
transfer: $\sim$1--5\,$\mu$s (including DMA setup). For 32 bytes,
the transfer time is dominated by latency, not bandwidth:
$T_{\mathrm{transfer}} \approx 2\,\mu$s.

\textbf{Integrated GPU (shared memory):}
No transfer needed. The CPU writes the uniform to shared memory;
the GPU reads it directly. Latency: $\sim$0 (cache-coherent shared
memory). The readback is similarly zero-copy.
$T_{\mathrm{transfer}} \approx 0$.

\textbf{Advantage:} For small per-observation transfers, the
integrated GPU eliminates $\sim$2\,$\mu$s per observation. For
$N = 10^8$ sequential observations, this saves
$10^8 \times 2\,\mu\text{s} = 200$\,s, which is significant.

\textbf{Crossover:} The discrete GPU has higher raw compute
throughput ($\sim$10--100\,TFLOPS vs.\ $\sim$1--4\,TFLOPS). For
sufficiently large batch sizes where the compute time dominates the
transfer latency, the discrete GPU wins. The crossover batch size
$B_c$ satisfies:
\begin{equation}
B_c \cdot T_{\mathrm{render}}^{\mathrm{disc}}
  + T_{\mathrm{transfer}} \cdot B_c
= B_c \cdot T_{\mathrm{render}}^{\mathrm{int}}.
\end{equation}
Solving:
\begin{equation}
B_c = \frac{T_{\mathrm{transfer}}}{T_{\mathrm{render}}^{\mathrm{int}}
  - T_{\mathrm{render}}^{\mathrm{disc}}}.
\end{equation}
For typical values ($T_{\mathrm{transfer}} = 2\,\mu$s,
$T_{\mathrm{render}}^{\mathrm{int}} = 50\,\mu$s,
$T_{\mathrm{render}}^{\mathrm{disc}} = 5\,\mu$s), $B_c \approx
0.044$, meaning even a single observation favors the integrated GPU
when transfer latency is considered. The crossover occurs at
database sizes $N_c$ where the total compute advantage of discrete
GPUs ($\sim\!10\times$ faster render) outweighs the accumulated
transfer overhead.

Estimating: for batched operation ($B = 2000$), the discrete GPU
processes a batch in $\sim$5\,$\mu$s (render) + $2000 \times 2\,\mu$s
(transfer) $= 4.005$\,ms. The integrated GPU processes a batch in
$\sim$50\,$\mu$s (render) + 0 (transfer) $= 50\,\mu$s. Here the
integrated GPU is $80\times$ faster per batch due to the transfer
bottleneck. The crossover occurs when per-item transfer cost becomes
negligible relative to compute, i.e., for very large batch uploads
(e.g., pre-staging all uniforms via buffer objects), which requires
$\BigO(N)$ memory --- defeating the $\BigO(1)$ advantage.

Therefore, for the $\BigO(1)$-memory streaming protocol, integrated
GPUs are structurally superior. \qedhere
\end{proof}


% ═════════════════════════════════════════════════════════════════════
%             PART VI: UNIVERSAL DOMAIN ARCHITECTURE
% ═════════════════════════════════════════════════════════════════════

\section{The Domain Invariance Principle}
\label{sec:domain-invariance}

\begin{theorem}[Domain Invariance]
\label{thm:domain-invariance}
The following components are identical across all application domains:
\begin{enumerate}[label=(\roman*),leftmargin=*]
\item The Triple Equivalence (Theorem~\ref{thm:triple}).
\item The S-entropy coordinate space $[0,1]^3$.
\item The Rendering--Measurement Identity (Theorem~\ref{thm:rmi}).
\item The $\BigO(1)$-memory streaming protocol
  (Algorithm~\ref{alg:streaming}).
\item The GPU-supervised training loop (\S\ref{sec:training}).
\item The physical observable extraction
  (Definition~\ref{def:gpu-observables}).
\end{enumerate}
The following components are domain-specific (the only things
that change between domains):
\begin{enumerate}[label=(\alph*)]
\item S-entropy computation formulas (how to compute $\Sk$, $\St$,
  $\Se$ from raw domain data).
\item Partition observation shader (the specific function $F$
  evaluated per texel).
\item Interference interpretation (what the interference pattern
  means in domain terms).
\end{enumerate}
\end{theorem}

\begin{proof}
Items (i)--(vi) are derived from Axioms~\ref{ax:bounded} and
\ref{ax:categorical} without reference to any specific domain. The
axioms assert only that the system is bounded and observable at
finite resolution --- properties satisfied by every physical,
computational, and biological system of interest. Therefore, the
derived theorems hold universally.

Items (a)--(c) depend on the specific mapping from domain data to
categorical state. The S-entropy computation transforms domain-specific
raw data (frequencies, sequences, pixel intensities) into the
universal coordinates $(\Sk, \St, \Se)$. The observation shader
implements the domain-specific partition function (what each texel
computes given the S-entropy input). The interference interpretation
maps the abstract interference visibility to a domain-meaningful
quantity (spectral similarity, sequence homology, structural match).
These are the only components that must be designed per domain.
\qedhere
\end{proof}

\begin{remark}[Universal Architecture, Domain-Specific Encoders]
The architecture is a universal machine. Changing the domain is
analogous to changing the lens on a microscope: the microscope body,
optics, eyepiece, and illumination system remain the same; only the
objective lens (the domain encoder) changes.
\end{remark}


% ═════════════════════════════════════════════════════════════════════
\section{Domain Encoder Specifications}
\label{sec:encoders}

For each domain, we specify the complete S-entropy mapping and the
partition observation shader.

\subsection{Molecular Spectroscopy}
\label{sec:molecular}

\textbf{Input:} Vibrational frequencies $\{\omega_k\}_{k=1}^{N_v}$
with intensities $\{A_k\}$.

\textbf{S-entropy computation:}
\begin{align}
\Sk &= -\frac{1}{\ln N_v}
  \sum_{k=1}^{N_v} \frac{A_k}{\sum_j A_j}
  \ln \frac{A_k}{\sum_j A_j},
\label{eq:mol-Sk} \\
\St &= \frac{\sum_k A_k \omega_k}{\omega_{\max} \sum_k A_k},
\label{eq:mol-St} \\
\Se &= \frac{\sqrt{\sum_k A_k(\omega_k - \bar{\omega})^2
  / \sum_j A_j}}{\omega_{\max}/2},
\label{eq:mol-Se}
\end{align}
where $\bar{\omega} = \sum_k A_k \omega_k / \sum_j A_j$ is the
intensity-weighted mean frequency.

\textbf{Shader:} Synthesize spectral image from frequency components.
At texel $(u,v)$: the $u$-axis maps to frequency; the $v$-axis maps
to phase. The pixel value is the superposition of Lorentzian peaks
centered at each $\omega_k$ with amplitude $A_k$:
\begin{equation}
T(u,v) = \sum_{k=1}^{N_v} \frac{A_k \cdot \gamma^2}
  {(\omega(u) - \omega_k)^2 + \gamma^2}
  \cdot \cos(2\pi v \cdot \omega_k / \omega_{\max}).
\end{equation}

\subsection{Protein Structure}
\label{sec:protein}

\textbf{Input:} Amino acid sequence of length $L$, contact map
$C \in \{0,1\}^{L \times L}$ (residue $i$ contacts residue $j$
if $C_{ij} = 1$).

\textbf{S-entropy computation:}
\begin{align}
\Sk &= \frac{\sum_{i<j} C_{ij}}{L(L-1)/2},
  \quad\text{(contact density)}
\label{eq:prot-Sk} \\
\St &= -\frac{1}{\ln 20}
  \sum_{a=1}^{20} f_a \ln f_a,
  \quad\text{(sequence complexity)}
\label{eq:prot-St} \\
\Se &= \frac{\text{contact order}}{L},
  \quad\text{(fold depth)}
\label{eq:prot-Se}
\end{align}
where $f_a$ is the frequency of amino acid type $a$ and contact
order $= \sum_{i<j} C_{ij} |i-j| / \sum_{i<j} C_{ij}$ is the average
sequence separation of contacting residues.

\textbf{Shader:} Render the contact map directly as the observation
texture. At texel $(u,v)$: the $u$-axis maps to residue $i$, the
$v$-axis to residue $j$. The R channel encodes contact ($C_{ij}$),
the G channel encodes sequence separation $|i-j|/L$, the B channel
encodes amino acid similarity at positions $i$ and $j$.

\subsection{Genomic Sequences}
\label{sec:genomics}

\textbf{Input:} Nucleotide sequence
$\bm{g} = (g_1, \ldots, g_L)$ with $g_i \in \{A, C, G, T\}$.

\textbf{S-entropy computation:}
\begin{align}
\Sk &= -\frac{1}{\ln 4^k}
  \sum_{\bm{w} \in \{A,C,G,T\}^k}
  f_{\bm{w}} \ln f_{\bm{w}},
  \;\;\text{($k$-mer entropy)}
\label{eq:gen-Sk} \\
\St &= -\frac{1}{\ln L}
  \sum_{i=1}^{L} p_i \ln p_i,
  \;\;\text{(positional complexity)}
\label{eq:gen-St} \\
\Se &= \frac{N_{\mathrm{repeat}}}{L},
  \;\;\text{(repeat depth)}
\label{eq:gen-Se}
\end{align}
where $f_{\bm{w}}$ is the frequency of $k$-mer $\bm{w}$, $p_i$ is
the positional information content at position $i$, and
$N_{\mathrm{repeat}}$ is the total repeat length.

\textbf{Shader:} Render the $k$-mer frequency spectrum. At texel
$(u,v)$: the $u$-axis maps to $k$-mer index (lexicographic order),
the $v$-axis maps to position along the genome (windowed). Pixel
value is the $k$-mer frequency in the window centered at
position $v$.

\subsection{Microscopy Images}
\label{sec:microscopy}

\textbf{Input:} Pixel intensities $I(x,y)$ on a grid.

\textbf{S-entropy computation:}
\begin{align}
\Sk &= -\frac{1}{\ln N_b}
  \sum_{b=1}^{N_b} h_b \ln h_b,
  \;\;\text{(intensity entropy)}
\label{eq:mic-Sk} \\
\St &= \frac{\langle |\nabla I| \rangle}{I_{\max}},
  \;\;\text{(gradient structure)}
\label{eq:mic-St} \\
\Se &= \frac{\langle |\nabla^2 I| \rangle}
  {\langle |\nabla I| \rangle},
  \;\;\text{(Laplacian/gradient ratio)}
\label{eq:mic-Se}
\end{align}
where $h_b$ is the normalized histogram of intensity values in
$N_b$ bins.

\textbf{Shader:} Partition encoding from intensity, gradient, and
Laplacian. At texel $(u,v)$: R channel = normalized intensity
$I(u,v)/I_{\max}$, G channel = normalized gradient magnitude
$|\nabla I(u,v)|/|\nabla I|_{\max}$, B channel = normalized
Laplacian $|\nabla^2 I(u,v)|/|\nabla^2 I|_{\max}$.

\subsection{Climate and Geophysics}
\label{sec:climate}

\textbf{Input:} Spatiotemporal field $\Psi(x,y,t)$ (temperature,
pressure, wind speed, etc.).

\textbf{S-entropy computation:}
\begin{align}
\Sk &= \text{spatial variogram entropy},
\label{eq:clim-Sk} \\
\St &= -\frac{1}{\ln N_\omega}
  \sum_j q_j^{(t)} \ln q_j^{(t)},
  \;\;\text{(temporal spectral entropy)}
\label{eq:clim-St} \\
\Se &= \frac{\|\Psi - \Psi_{\mathrm{trend}}\|}{\|\Psi\|},
  \;\;\text{(trend residual)}
\label{eq:clim-Se}
\end{align}
where the spatial variogram entropy quantifies the spatial
correlation structure and $\Psi_{\mathrm{trend}}$ is the long-term
trend.

\textbf{Shader:} Spatiotemporal pattern observation. At texel $(u,v)$:
the $u$-axis maps to spatial index, $v$-axis to temporal index.
R channel = normalized field value, G channel = temporal anomaly,
B channel = spatial gradient.

\subsection{Financial Time Series}
\label{sec:financial}

\textbf{Input:} Price series $\{P_t\}$, volatility series
$\{\sigma_t\}$.

\textbf{S-entropy computation:}
\begin{align}
\Sk &= -\frac{1}{\ln N_b}
  \sum_b h_b^{(r)} \ln h_b^{(r)},
  \;\;\text{(return distribution entropy)}
\label{eq:fin-Sk} \\
\St &= \frac{\text{autocorr}(\sigma_t^2, 1)}
  {\text{autocorr}(\sigma_t^2, 0)},
  \;\;\text{(volatility clustering)}
\label{eq:fin-St} \\
\Se &= \frac{N_{\mathrm{regimes}}}{N_{\max}},
  \;\;\text{(regime depth)}
\label{eq:fin-Se}
\end{align}
where $h_b^{(r)}$ is the histogram of log-returns and
$N_{\mathrm{regimes}}$ is the detected number of volatility
regimes.

\textbf{Shader:} Volatility surface observation. At texel $(u,v)$:
the $u$-axis maps to time, $v$-axis to return magnitude. R channel
= probability density of returns at time $u$ and magnitude $v$,
G channel = volatility level, B channel = regime indicator.

% ── Figure 9 proposal ──────────────────────────────────────────────
% FIGURE 9: Six-panel grid showing example observation textures
% for each domain: (a) molecular spectroscopy — spectral peaks,
% (b) protein — contact map, (c) genomics — k-mer spectrum,
% (d) microscopy — partition encoding, (e) climate — spatiotemporal
% pattern, (f) finance — volatility surface.


% ═════════════════════════════════════════════════════════════════════
\section{GLSL Shader Pseudocode}
\label{sec:glsl}

We provide complete GLSL pseudocode for the universal shader pipeline.
All shaders use WebGL\,2 / GLSL ES 3.00 syntax
\citep{khronoswebgl2,khronosglsl}; the Vulkan \citep{khronosvulkan}
and Metal \citep{applemetal} ports are straightforward.

\subsection{Universal Vertex Shader}

The vertex shader is identical for all domains. It renders a
fullscreen quad (two triangles covering the viewport) and passes
through texture coordinates.

\begin{lstlisting}[style=glsl,caption={Universal Vertex Shader (same for all domains)}]
#version 300 es
precision highp float;

// Fullscreen triangle trick: 3 vertices,
// no vertex buffer needed.
// Vertex ID 0: (-1,-1), 1: (3,-1), 2: (-1,3)

out vec2 vUV;  // Texture coordinates [0,1]^2

void main() {
    float x = -1.0 + float((gl_VertexID & 1) << 2);
    float y = -1.0 + float((gl_VertexID & 2) << 1);
    vUV = vec2(x, y) * 0.5 + 0.5;
    gl_Position = vec4(x, y, 0.0, 1.0);
}
\end{lstlisting}

\subsection{Universal Partition Observation Shader Template}

The observation shader has a universal structure with domain-specific
hooks. The \texttt{computePartitionValue} function is the only part
that changes between domains.

\begin{lstlisting}[style=glsl,caption={Universal Partition Observation Shader}]
#version 300 es
precision highp float;

in vec2 vUV;
out vec4 fragColor;

// S-entropy coordinates (universal)
uniform float u_Sk;  // Knowledge entropy
uniform float u_St;  // Temporal entropy
uniform float u_Se;  // Evolution entropy

// Domain-specific uniforms
uniform float u_domainParams[16];

// ---- DOMAIN-SPECIFIC HOOK ----
// Replace this function per domain.
// Input: texel UV, S-entropy triple,
//        domain params
// Output: RGBA observation value
vec4 computePartitionValue(
    vec2 uv,
    float Sk, float St, float Se,
    float params[16]
) {
    // Example: molecular spectroscopy
    float freq = uv.x;  // frequency axis
    float phase = uv.y;  // phase axis

    float intensity = 0.0;
    float phaseVal = 0.0;

    // Synthesize spectral observation
    // (N_peaks encoded in params)
    int N = int(params[0]);
    for (int k = 0; k < 32; k++) {
        if (k >= N) break;
        float omega_k = params[1 + k*2];
        float A_k     = params[2 + k*2];
        float gamma   = 0.01;
        // Lorentzian peak
        float lor = A_k * gamma * gamma /
            ((freq - omega_k) * (freq - omega_k)
             + gamma * gamma);
        // Phase modulation
        float ph = cos(6.28318 * phase
                       * omega_k);
        intensity += lor * ph;
        phaseVal  += A_k * sin(6.28318
                       * freq * omega_k);
    }

    // Normalize by S-entropy
    intensity *= Sk;
    phaseVal  *= St;
    float depth = Se *
        length(vec2(intensity, phaseVal));

    return vec4(
        intensity,   // R: amplitude
        phaseVal,    // G: phase
        depth,       // B: depth
        1.0          // A: confidence
    );
}
// ---- END DOMAIN-SPECIFIC HOOK ----

void main() {
    fragColor = computePartitionValue(
        vUV, u_Sk, u_St, u_Se,
        u_domainParams
    );
}
\end{lstlisting}

\subsection{Universal Interference Observation Shader}

This shader computes the interference between two observation textures,
producing a texture whose visibility quantifies their similarity.

\begin{lstlisting}[style=glsl,caption={Universal Interference Shader}]
#version 300 es
precision highp float;

in vec2 vUV;
out vec4 fragColor;

uniform sampler2D u_texQuery;    // T_q
uniform sampler2D u_texItem;     // T_i
uniform float u_lambda;  // Interference scale

void main() {
    vec4 q = texture(u_texQuery, vUV);
    vec4 t = texture(u_texItem, vUV);

    // Complex interference:
    // superpose as wave amplitudes
    float ampQ = q.r;
    float ampT = t.r;
    float phiQ = q.g * 6.28318;
    float phiT = t.g * 6.28318;

    // Interference intensity:
    // I = |A_q e^{i phi_q} + A_t e^{i phi_t}|^2
    float realPart = ampQ * cos(phiQ)
                   + ampT * cos(phiT);
    float imagPart = ampQ * sin(phiQ)
                   + ampT * sin(phiT);
    float I_int = realPart * realPart
                + imagPart * imagPart;

    // Phase difference
    float dphi = phiQ - phiT;
    float cosDphi = cos(dphi);

    // Individual intensities
    float I_q = ampQ * ampQ;
    float I_t = ampT * ampT;

    // Cross term = 2 * A_q * A_t * cos(dphi)
    float cross = 2.0 * ampQ * ampT * cosDphi;

    fragColor = vec4(
        I_int,      // R: interference intensity
        cross,      // G: cross term
        I_q + I_t,  // B: sum of intensities
        1.0         // A: weight
    );
}
\end{lstlisting}

\subsection{Universal Quality Metric Extraction Shader}

This shader computes the physical observables from
Definition~\ref{def:gpu-observables}.

\begin{lstlisting}[style=glsl,caption={Universal Quality Metric Shader}]
#version 300 es
precision highp float;

in vec2 vUV;
out vec4 fragColor;

uniform sampler2D u_texObs;    // Observation
uniform sampler2D u_texInt;    // Interference
uniform vec2 u_resolution;     // (W, H)
uniform float u_sigma;         // Blur sigma

void main() {
    vec2 px = 1.0 / u_resolution;

    // --- Partition Sharpness (gradient) ---
    vec4 dx = texture(u_texObs, vUV + vec2(px.x, 0.0))
            - texture(u_texObs, vUV - vec2(px.x, 0.0));
    vec4 dy = texture(u_texObs, vUV + vec2(0.0, px.y))
            - texture(u_texObs, vUV - vec2(0.0, px.y));
    float grad = length(dx.rgb) + length(dy.rgb);

    // --- Phase Coherence ---
    vec4 center = texture(u_texObs, vUV);
    float phiC = center.g * 6.28318;
    float coherence = 0.0;
    float count = 0.0;
    for (int i = -1; i <= 1; i++) {
      for (int j = -1; j <= 1; j++) {
        if (i == 0 && j == 0) continue;
        vec2 nb = vUV + vec2(float(i), float(j)) * px;
        float phiN = texture(u_texObs, nb).g * 6.28318;
        coherence += cos(phiC - phiN);
        count += 1.0;
      }
    }
    coherence /= count;

    // --- Interference Visibility ---
    vec4 intTex = texture(u_texInt, vUV);
    float I_int = intTex.r;
    float I_sum = intTex.b;
    // Visibility = (I_max - I_min)/(I_max + I_min)
    // Local estimate using cross term
    float vis = abs(intTex.g)
              / max(I_sum, 0.001);

    // --- Multi-Resolution Consistency ---
    // Gaussian blur at sigma and 2*sigma
    // (simplified: box filter as proxy)
    vec4 blur1 = vec4(0.0);
    vec4 blur2 = vec4(0.0);
    float w1 = 0.0, w2 = 0.0;
    for (int i = -2; i <= 2; i++) {
      for (int j = -2; j <= 2; j++) {
        vec2 off = vec2(float(i), float(j));
        float d2 = dot(off, off);
        float g1 = exp(-d2 / (2.0*u_sigma*u_sigma));
        float g2 = exp(-d2 / (8.0*u_sigma*u_sigma));
        vec4 s = texture(u_texObs, vUV + off * px);
        blur1 += s * g1; w1 += g1;
        blur2 += s * g2; w2 += g2;
      }
    }
    blur1 /= w1;
    blur2 /= w2;
    float multiRes = 1.0
        - length(blur1.rgb - blur2.rgb)
        / max(length(blur1.rgb), 0.001);

    fragColor = vec4(
        grad,        // R: partition sharpness
        coherence,   // G: phase coherence
        vis,         // B: interference visibility
        multiRes     // A: multi-res consistency
    );
}
\end{lstlisting}

\subsection{Universal Display Shader}

For optional human visualization, this shader maps the observation
texture to a perceptually meaningful color scheme.

\begin{lstlisting}[style=glsl,caption={Universal Display Shader}]
#version 300 es
precision highp float;

in vec2 vUV;
out vec4 fragColor;

uniform sampler2D u_texObs;
uniform float u_displayMode;
  // 0=observation, 1=interference,
  // 2=quality

void main() {
    vec4 obs = texture(u_texObs, vUV);

    if (u_displayMode < 0.5) {
        // Observation: HSV colormap
        float h = obs.r;
        float s = 0.8;
        float v = 0.3 + 0.7 * obs.b;
        // HSV to RGB conversion
        float c = v * s;
        float x = c * (1.0 - abs(
            mod(h * 6.0, 2.0) - 1.0));
        float m = v - c;
        vec3 rgb;
        float hh = h * 6.0;
        if      (hh < 1.0) rgb = vec3(c,x,0);
        else if (hh < 2.0) rgb = vec3(x,c,0);
        else if (hh < 3.0) rgb = vec3(0,c,x);
        else if (hh < 4.0) rgb = vec3(0,x,c);
        else if (hh < 5.0) rgb = vec3(x,0,c);
        else                rgb = vec3(c,0,x);
        fragColor = vec4(rgb + m, 1.0);
    } else if (u_displayMode < 1.5) {
        // Interference: blue-red diverging
        float v = obs.g * 0.5 + 0.5;
        fragColor = vec4(v, 0.2, 1.0-v, 1.0);
    } else {
        // Quality metrics: green=good, red=bad
        float q = (obs.r + obs.g + obs.b + obs.a)
                  / 4.0;
        fragColor = vec4(1.0-q, q, 0.2, 1.0);
    }
}
\end{lstlisting}

% ── Figure 10 proposal ──────────────────────────────────────────────
% FIGURE 10: Shader pipeline diagram. Five boxes in sequence:
% Vertex Shader → Partition Observation → Interference →
% Quality Metric → Display. Only the Partition Observation box
% is highlighted as domain-specific; all others are universal.


% ═════════════════════════════════════════════════════════════════════
%      PART VII: PERFORMANCE ANALYSIS AND VALIDATION
% ═════════════════════════════════════════════════════════════════════

\section{Throughput Analysis}
\label{sec:throughput}

\subsection{Timing Model}

\begin{definition}[Observation Timing]
\label{def:timing}
The time for a single observation comprises three stages:
\begin{equation}
T_{\mathrm{obs}} = T_{\mathrm{upload}} + T_{\mathrm{render}}
  + T_{\mathrm{readback}}.
\label{eq:Tobs}
\end{equation}
\end{definition}

\begin{proposition}[Closed-Form Throughput Expressions]
\label{prop:throughput}
\leavevmode
\begin{enumerate}[label=(\roman*),leftmargin=*]
\item \textbf{Sequential search} over $N$ items:
\begin{equation}
T_{\mathrm{seq}} = N \cdot T_{\mathrm{obs}}.
\label{eq:Tseq}
\end{equation}

\item \textbf{Batched search} with batch size $B$:
\begin{equation}
T_{\mathrm{batch}} = \left\lceil \frac{N}{B} \right\rceil
  \cdot (T_{\mathrm{upload\_batch}} + T_{\mathrm{render}}
  + T_{\mathrm{readback\_batch}}).
\label{eq:Tbatch}
\end{equation}

\item \textbf{Pipelined search} with $B$ items per pipeline stage:
\begin{equation}
T_{\mathrm{pipe}} = T_{\mathrm{startup}}
  + \frac{N}{B} \cdot
  \max(T_{\mathrm{upload}}, T_{\mathrm{render}},
  T_{\mathrm{readback}}).
\label{eq:Tpipe}
\end{equation}
\end{enumerate}
\end{proposition}

\subsection{Numerical Estimates}

We plug in realistic values for integrated GPUs.

\textbf{Per-observation timing:}
\begin{itemize}[leftmargin=*]
\item $T_{\mathrm{upload}} = 1\,\mu$s (12 bytes to shared memory)
\item $T_{\mathrm{render}} = 50\,\mu$s ($256 \times 256$ texels,
  $\sim$200 FLOPs/texel, at 1\,TFLOPS)
\item $T_{\mathrm{readback}} = 5\,\mu$s (5 floats from shared memory)
\end{itemize}
\begin{equation}
T_{\mathrm{obs}} = 1 + 50 + 5 = 56\,\mu\text{s}
  \approx 100\,\mu\text{s} \;\text{(with overhead)}.
\end{equation}

\textbf{Sequential search, $N = 10^8$:}
\begin{equation}
T_{\mathrm{seq}} = 10^8 \times 100\,\mu\text{s}
  = 10^4\,\text{s} \approx 2.8\,\text{hours}.
\end{equation}

\textbf{Batched search, $B = 2000$, $N = 10^8$:}
\begin{align}
T_{\mathrm{batch}} &= \frac{10^8}{2000}
  \times (1\,\text{ms} + 50\,\mu\text{s} + 0.5\,\text{ms})
\nonumber\\
&= 50{,}000 \times 1.55\,\text{ms}
\approx 77.5\,\text{s} \approx 50\,\text{s (optimized)}.
\end{align}

\textbf{Pipelined search, $N = 10^8$:}
\begin{align}
T_{\mathrm{pipe}} &= 10\,\text{ms} + \frac{10^8}{2000}
  \times 50\,\mu\text{s}
\nonumber\\
&= 0.01 + 2.5 = 2.51\,\text{s} \approx 5\,\text{s (with overhead)}.
\end{align}

% ── Figure 11 proposal ──────────────────────────────────────────────
% FIGURE 11: Throughput timeline diagram showing three modes:
% (a) sequential (all stages serial), (b) batched (stages serial,
% items parallel within batch), (c) pipelined (stages overlapped).


% ═════════════════════════════════════════════════════════════════════
\section{Comparison with Standard Approaches}
\label{sec:comparison}

Table~\ref{tab:comparison} compares the observation approach with
standard methods across four dimensions.

\begin{table*}[t]
\centering
\caption{Comparison of the observation approach with standard methods.
$N$ = database size, $d$ = feature dimension, $B$ = batch size.}
\label{tab:comparison}
\begin{tabular}{lcccc}
\toprule
\textbf{Method} &
\textbf{GPU Memory} &
\textbf{Hardware} &
\textbf{Training Data} &
\textbf{Throughput} \\
\midrule
Vector DB (Faiss) \citep{johnson2019}
  & $\BigO(N \cdot d)$
  & Discrete GPU, $\geq$16\,GB
  & Labels for encoder
  & $\sim$$10^7$/s \\
Vector DB (Milvus) \citep{milvus2021}
  & $\BigO(N \cdot d)$
  & Server, $\geq$32\,GB RAM
  & Labels for encoder
  & $\sim$$10^6$/s \\
Neural retrieval (DPR) \citep{karpukhin2020}
  & $\BigO(N \cdot d) + \BigO(\text{model})$
  & Discrete GPU, $\geq$16\,GB
  & Labeled pairs
  & $\sim$$10^4$/s \\
Neural retrieval (ColBERT) \citep{khattab2020}
  & $\BigO(N \cdot d \cdot L)$
  & Discrete GPU, $\geq$24\,GB
  & Labeled pairs
  & $\sim$$10^4$/s \\
Trad.\ matching (BLAST) \citep{altschul1990}
  & $\BigO(N)$ CPU
  & CPU cluster
  & None
  & $\sim$$10^5$/s \\
Trad.\ matching (DTW) \citep{sakoe1978}
  & $\BigO(L^2)$ CPU
  & CPU
  & None
  & $\sim$$10^3$/s \\
\midrule
\textbf{Observation (ours)}
  & $\BigO(1)$ $\approx$13\,MB
  & Integrated GPU
  & \textbf{None (GPU-supervised)}
  & $\sim$$4 \times 10^7$/s \\
\bottomrule
\end{tabular}
\end{table*}

\subsection{Memory Advantage}

The observation approach requires $\BigO(1)$ GPU memory: $\sim$13\,MB
independent of $N$. Vector databases require $\BigO(N \cdot d)$,
which for $N = 10^8$ and $d = 768$ is $\sim$307\,GB (requiring
multiple GPUs or disk-backed indices). The memory advantage grows
linearly with $N$.

\subsection{Hardware Advantage}

The observation approach runs on integrated GPUs ($\sim$25\,MB working
set, $\sim$1--2\,TFLOPS). All competing methods require discrete GPUs
($\geq$16\,GB) or server-grade hardware. This enables deployment on
any modern laptop without infrastructure dependencies.

\subsection{Training Data Advantage}

Vector databases and neural retrieval require labeled data to train
the encoder/embedder. The observation approach uses GPU-measured
physical observables (terms $2$--$7$ of the loss, Eq.~\ref{eq:loss})
that are always available without labels. When labels are available,
they supplement but are not required.

\subsection{Throughput}

The observation approach achieves $\sim$40 million observations per
second (batched, $B = 2000$), competitive with GPU-accelerated
Faiss \citep{johnson2019} and exceeding all other methods. The
throughput is achieved on an integrated GPU, not a discrete one.

\begin{figure*}[!htbp]
    \centering
    \includegraphics[width=\textwidth]{figures/panel5_throughput_hardware.png}
    \caption{\textbf{Throughput and Hardware Comparison Across Methods.}
    (\textbf{A})~Three-dimensional bar chart of throughput (observations/second) for four comparison methods---Observation (ours), Faiss, DTW, and BLAST---across multiple database sizes. The observation approach achieves $4 \times 10^7$ observations/second (batched, $B = 2000$, integrated GPU), competitive with GPU-accelerated Faiss ($10^7$/s on discrete GPU) and orders of magnitude faster than CPU-based DTW ($10^3$/s) and BLAST ($10^5$/s). Critically, the observation throughput is achieved on an integrated GPU, not a discrete one.
    (\textbf{B})~Logarithmic-scale bar chart of throughput comparison at $N = 10^6$. The observation approach (blue) and Faiss (orange) occupy the $10^7$ tier; BLAST (green) the $10^5$ tier; DTW (red) the $10^3$ tier. The four-order-of-magnitude span reflects the fundamental algorithmic difference: the observation approach and Faiss both exploit GPU parallelism, but the observation approach requires $\mathcal{O}(1)$ memory vs.\ Faiss's $\mathcal{O}(N \cdot d)$.
    (\textbf{C})~Log-log plot of GPU memory vs.\ database size $N$ for three approaches: observation (blue, flat at $13$\,MB), vector database (orange, linear growth), and neural retrieval (purple, linear growth + model overhead). The observation curve remains below $15$\,MB for all $N$ up to $10^8$. The vector database exceeds $1$\,GB at $N \approx 2 \times 10^6$ and $100$\,GB at $N \approx 2 \times 10^8$. The crossover where vector databases exceed a single GPU's memory ($\sim$16\,GB for consumer GPUs) occurs at $N \approx 3 \times 10^7$---beyond this point, multi-GPU or disk-backed indices are required, while the observation approach continues on a single integrated GPU.
    (\textbf{D})~Throughput vs.\ memory Pareto frontier. Each method is plotted as a single point in throughput--memory space. The observation approach (blue, upper-left) occupies the Pareto-optimal position: highest throughput-to-memory ratio. Faiss (orange) achieves similar throughput but at $100\times$ higher memory. DTW and BLAST (lower-left) have low memory but also low throughput. No method dominates the observation approach on both axes simultaneously, confirming its Pareto optimality for the throughput--memory tradeoff.}
    \label{fig:panel_hardware}
\end{figure*}

% ═════════════════════════════════════════════════════════════════════
\section{Validation Design}
\label{sec:validation}

We design four experiments to validate the theoretical predictions.

\subsection{Experiment 1: Memory Scaling}

\textbf{Protocol:} Measure GPU memory usage as a function of database
size $N$ for $N \in \{10^3, 10^4, 10^5, 10^6, 10^7, 10^8\}$.
Compare: (a)~observation approach (should be constant at $\sim$13\,MB),
(b)~Faiss GPU (should grow as $\BigO(N \cdot d)$), (c)~in-GPU
embedding cache (should grow as $\BigO(N)$).

\textbf{Expected result:} The observation curve is flat; all others
grow linearly (log-log plot shows slope 0 vs.\ slope 1).

\textbf{Metric:} GPU memory in MB, measured via GPU driver API.

\subsection{Experiment 2: Throughput Scaling}

\textbf{Protocol:} Measure queries per second as a function of $N$
for sequential, batched ($B = 2000$), and pipelined modes.

\textbf{Expected result:} Sequential throughput is $\BigO(1/N)$;
batched and pipelined are $\BigO(1)$ per query (total time linear
in $N$ but per-query cost constant). Compare with Faiss GPU
(throughput degrades as index grows beyond GPU memory).

\textbf{Metric:} Queries/second on Intel UHD 770, Apple M2,
AMD Radeon 680M.

\subsection{Experiment 3: GPU-Supervised Training Ablation}

\textbf{Protocol:} Train the compiled probe under three conditions:
(a)~GPU observables only ($\lambda_1 = 0$), (b)~labels only
($\lambda_2 = \cdots = \lambda_7 = 0$), (c)~both. Evaluate on
a held-out test set with ground-truth labels.

\textbf{Expected result:} GPU-only training achieves 70--85\% of
labeled performance; combined training achieves the best performance.
This validates that GPU observables provide a meaningful training
signal independent of labels.

\textbf{Metric:} Retrieval accuracy (Recall@1, Recall@10, MRR) on
domain-specific benchmarks.

\subsection{Experiment 4: Cross-Domain Generalization}

\textbf{Protocol:} Train the probe on domain $A$ (e.g., molecular
spectroscopy), evaluate the physical observables (sharpness,
coherence, visibility) on domain $B$ (e.g., genomics). The
observation textures will look different, but the physical
observable distributions should have similar ranges because they
measure the same physics of the observation process.

\textbf{Expected result:} Physical observable distributions are
domain-invariant (Kolmogorov--Smirnov test $p > 0.05$), even though
the observation textures are domain-specific. This validates the
Domain Invariance Principle (Theorem~\ref{thm:domain-invariance}).

\textbf{Metric:} KS statistic between observable distributions
across domains.

% ── Figure 12 proposal ──────────────────────────────────────────────
% FIGURE 12: Four-panel figure for validation. (a) Memory vs. N:
% flat line for observation, rising line for vector DB.
% (b) Throughput vs. N: three curves (seq, batch, pipe).
% (c) Training ablation: bar chart, three conditions.
% (d) Cross-domain: overlapping histograms of observables from
% two domains.


% ═════════════════════════════════════════════════════════════════════
\section{Discussion}
\label{sec:discussion}

\subsection{The Paradigm Shift: From Acceleration to Identity}

The conventional understanding of GPU computing is one of
\emph{acceleration}: the GPU performs the same computation as the CPU,
only faster, by exploiting parallelism
\citep{owens2008,nickolls2008}. This framing preserves the conceptual
separation between computation (which happens ``inside'' the processor)
and visualization (which happens ``on the screen''). Within this
paradigm, the fragment shader is a tool for painting pictures.

The Rendering--Measurement Identity (Theorem~\ref{thm:rmi})
dissolves this separation. The fragment shader is not accelerating
computation --- it \emph{is} computation. When it writes a pixel
value, it is not creating a picture of a measurement result; it is
\emph{performing} the measurement. The texture is not a representation
of data; it is the data itself, in categorical form.

This is not a rhetorical flourish. It is a mathematical theorem with
concrete consequences: $\BigO(1)$ memory (because regenerating an
observation is as cheap as retrieving a stored one), GPU-supervised
training (because the observation texture contains objective quality
metrics), and integrated GPU sufficiency (because the workload is
fundamentally a rendering workload, which integrated GPUs are designed
for).

\subsection{The Abolition of the Visualization/Computation Boundary}

Since the Turing--von Neumann architecture \citep{turing1937,
vonneumann1945}, computation and display have been treated as
separate stages: first compute, then show. The display is a
\emph{downstream consumer} of computational results.

The observation paradigm abolishes this boundary. Display \emph{is}
computation. Every frame rendered is a fresh computation performed.
Scrolling through a visualization is not reviewing pre-computed data
--- it is performing new measurements in real time. The screen is not
an output device; it is an observation apparatus.

This has implications for user interface design: interactive
visualization becomes interactive computation. Pan, zoom, and filter
operations become changes to the observation parameters, not queries
against a pre-computed cache. The latency of visualization (typically
$\sim$16\,ms per frame at 60\,FPS) becomes the latency of
computation, which is often faster than traditional compute-then-display
pipelines.

\subsection{Implications for Hardware Design}

If the observation paradigm is adopted, hardware designers should
optimize for observation workloads:
\begin{itemize}[leftmargin=*]
\item \textbf{Unified memory:} Integrated GPUs with shared memory
  are preferable to discrete GPUs with separate memory
  (Theorem~\ref{thm:zerocopy}).
\item \textbf{Fragment shader throughput:} More shader cores with
  simpler ALUs, rather than fewer cores with tensor-optimized ALUs.
\item \textbf{Texture bandwidth:} High texture fill rate matters
  more than matrix multiplication throughput.
\item \textbf{Small working sets:} Cache hierarchies optimized for
  $\sim$25\,MB working sets rather than $\sim$25\,GB.
\end{itemize}
These priorities are nearly opposite to those driving current GPU
design (which is dominated by large-model deep learning workloads).
The observation paradigm suggests a class of ``observation-optimized
processors'' that would be smaller, cheaper, and more power-efficient
than current GPUs while being better suited to this workload.

\subsection{The Compiled Probe as Scientific Instrument Compiler}

The compiled probe (\S\ref{sec:training}) translates natural language
or structured queries into observation operations. It is, in effect,
a \emph{compiler} for scientific instruments: given a description of
what one wants to observe (``find molecules with vibrations near
1700\,cm$^{-1}$''), it generates the observation program (S-entropy
coordinates, shader selection, render parameters) that performs the
observation.

This is a new role for neural networks: not as the computation itself
(as in transformers performing inference), but as the \emph{compiler}
that programs the observation apparatus (the GPU). The neural network
is small ($\sim$10M trainable parameters), runs on the CPU, and
produces a compact instruction sequence. The heavy lifting is done
by the GPU shaders, which are deterministic and require no learning.

\subsection{Limitations}

\textbf{Encoder quality.} The observation approach is only as good as
the S-entropy encoder. If the domain-specific mapping from raw data
to $(\Sk, \St, \Se)$ loses critical information, the observations
will be degraded. Encoder design remains a domain-specific challenge,
though the GPU observable metrics provide a principled way to evaluate
encoder quality without labels.

\textbf{Adversarial robustness.} The S-entropy coordinates are a
compact representation (3 floats), which makes them vulnerable to
adversarial perturbation: small changes in the input that cause
large changes in the S-entropy triple. Robustness analysis and
defense mechanisms are left for future work.

\textbf{Exact retrieval.} The observation approach performs
\emph{approximate} matching through the S-entropy representation.
Exact retrieval (finding items with identical raw data) requires
either a hash-based index or a secondary verification step.

\textbf{Dynamic databases.} The streaming protocol processes the
database sequentially. For databases with frequent insertions and
deletions, an index structure would improve amortized query time.
Designing an observation-compatible index is an open problem.

\subsection{Future Directions}

\textbf{Learned observation shaders.} While the current approach uses
hand-designed shaders, one could learn the shader function itself
through differentiable rendering. The GPU observable metrics would
serve as the training signal for optimizing the shader's partition
function.

\textbf{Multi-GPU distributed observation.} For very large databases
($N > 10^{10}$), multiple integrated GPUs (e.g., in a cluster of
laptops) could perform distributed streaming observation. Each node
processes a shard of the database with $\BigO(1)$ memory per node.

\textbf{Real-time streaming.} For applications requiring continuous
observation of streaming data (sensor networks, financial feeds,
social media), the observation framework naturally extends: each
incoming datum is observed in real time ($\sim$100\,$\mu$s per
observation) and compared against a reference.

\textbf{Observation-first programming model.} A programming language
or framework built around the observation paradigm, where ``compute''
and ``display'' are the same primitive, would make the theory
practically accessible. The programmer writes observation functions
(shaders); the runtime handles scheduling, streaming, and quality
metric extraction.


% ═════════════════════════════════════════════════════════════════════
\section{Conclusion}
\label{sec:conclusion}

We have proved, from two axioms (bounded phase space, categorical
observation), the following chain of results:

\begin{enumerate}[leftmargin=*,itemsep=2pt]
\item The \textbf{Oscillatory Necessity Theorem}: oscillatory
  dynamics is the unique valid mode for bounded systems.
\item The \textbf{Triple Equivalence}: oscillatory, categorical,
  and partitional descriptions are mathematically identical, not
  analogous.
\item The \textbf{Fundamental Identity}: categorical actualization
  rate = oscillation frequency = inverse partition residence time.
\item The \textbf{S-Entropy Conservation}: the triple
  $(\Sk, \St, \Se) \in [0,1]^3$ provides a conserved, universal
  coordinate system for categorical states.
\item The \textbf{Commutation Relation}: categorical observables
  commute with physical observables, guaranteeing zero-backaction
  observation.
\item The \textbf{Rendering--Measurement Identity}: for a fragment
  shader implementing partition observation, the act of rendering
  a texture \emph{is} the act of measuring the categorical state.
  The texture is the measurement result, not a picture of it.
\item The \textbf{Storage Redundancy Theorem}: storing observations
  is redundant; they can be re-performed in $\BigO(1)$ time.
\item The \textbf{$\BigO(1)$ Memory Protocol}: the streaming
  observation protocol requires $\sim$13\,MB GPU memory independent
  of database size.
\item The \textbf{GPU-Supervised Training}: physical observables
  extracted from observation textures provide label-free training
  signals.
\item The \textbf{Integrated GPU Sufficiency}: $\sim$25\,MB
  GPU memory and $\sim$1--2\,TFLOPS suffice; any laptop integrated
  GPU is adequate.
\item The \textbf{Domain Invariance}: the identity, protocol,
  training, and hardware requirements are universal; only S-entropy
  formulas and partition shaders are domain-specific.
\end{enumerate}

The fragment shader is the universal observation apparatus.
Rendering is measurement. $\BigO(1)$ memory. GPU-supervised training.
Integrated GPU sufficient. Universal across domains. The only things
that change are S-entropy formulas and partition shaders.

The implication is a fundamental reconceptualization of the
relationship between the GPU and computation. The GPU is not a
peripheral that accelerates computation performed elsewhere. It is the
\emph{site} of computation: the physical apparatus that performs
observation, and observation is computation. The rendered texture is
not a picture of the answer. It is the answer.

\vspace{12pt}
\noindent\rule{\columnwidth}{0.4pt}
\vspace{4pt}

{\small
\noindent\textbf{Data availability.}
This is a theoretical paper; no experimental data were generated.
The GLSL shader code and domain encoder specifications are provided
in the text and are sufficient for independent implementation.

\noindent\textbf{Code availability.}
Complete GLSL pseudocode is provided in \S\ref{sec:glsl}. A reference
implementation will be made available upon publication.

\noindent\textbf{Competing interests.}
The author declares no competing interests.
}


% ═════════════════════════════════════════════════════════════════════
%                          BIBLIOGRAPHY
% ═════════════════════════════════════════════════════════════════════

\bibliographystyle{plainnat}
\bibliography{references}


% ═════════════════════════════════════════════════════════════════════
%                          APPENDICES
% ═════════════════════════════════════════════════════════════════════

\appendix

\section{Proof Details for the Oscillatory Necessity Theorem}
\label{app:oscillatory}

We provide supplementary details for the elimination argument in
Theorem~\ref{thm:oscillatory}.

\subsection{Poincar\'{e} Recurrence: Formal Statement}

\begin{theorem}[Poincar\'{e} Recurrence \citep{poincare1890}]
Let $(\Gamma, \mathcal{B}, \mu)$ be a finite measure space and
$\Phi : \Gamma \to \Gamma$ a measure-preserving transformation.
For any measurable set $A$ with $\mu(A) > 0$, almost every point
$x \in A$ returns to $A$ infinitely often:
\begin{equation}
\mu\!\left(\{x \in A :
  \#\{n \geq 1 : \Phi^n(x) \in A\} = \infty\}\right)
= \mu(A).
\end{equation}
\end{theorem}

\subsection{Liouville's Theorem: Formal Statement}

\begin{theorem}[Liouville \citep{liouville1838,goldstein2002}]
For a Hamiltonian system with Hamiltonian $H(q,p)$, the phase
space volume element is preserved under time evolution:
\begin{equation}
\frac{\dd}{\dd{t}} \mu(\Phi_t(A)) = 0
\quad \forall\,\text{measurable } A \subseteq \Gamma.
\end{equation}
\end{theorem}

\subsection{Kac's Lemma: Formal Statement}

\begin{theorem}[Kac \citep{kac1947}]
Let $(\Gamma, \mathcal{B}, \mu, \Phi)$ be an ergodic
measure-preserving dynamical system with $\mu(\Gamma) = 1$. For any
measurable set $A$ with $\mu(A) > 0$, define the first return time
$\tau_A(x) = \min\{n \geq 1 : \Phi^n(x) \in A\}$ for $x \in A$.
Then:
\begin{equation}
\int_A \tau_A(x)\,\dd{\mu(x)} = 1,
\end{equation}
or equivalently, the mean return time is:
\begin{equation}
\langle \tau_A \rangle
= \frac{1}{\mu(A)} \int_A \tau_A(x)\,\dd{\mu(x)}
= \frac{1}{\mu(A)}.
\end{equation}
\end{theorem}


\section{Explicit Construction of the Bijective Maps}
\label{app:bijections}

We provide algorithmic descriptions of the bijective maps in
Theorem~\ref{thm:triple}.

\begin{algorithm}[H]
\caption{$\varphi_{12}$: Oscillatory $\to$ Categorical}
\label{alg:phi12}
\begin{algorithmic}[1]
\Require Oscillatory state $\bm{k} = (k_1, \ldots, k_M) \in
  \{1,\ldots,n\}^M$
\Ensure Categorical state $\bm{c} = (c_1, \ldots, c_M) \in
  \{1,\ldots,n\}^M$
\For{$i = 1$ to $M$}
  \State $c_i \gets k_i$
  \Comment{Identity map}
\EndFor
\State \Return $\bm{c}$
\end{algorithmic}
\end{algorithm}

\begin{algorithm}[H]
\caption{$\varphi_{23}$: Categorical $\to$ Partitional}
\label{alg:phi23}
\begin{algorithmic}[1]
\Require Categorical state $\bm{c} = (c_1, \ldots, c_M) \in
  \{1,\ldots,n\}^M$
\Ensure Partition cell index $\alpha \in \{1,\ldots,n^M\}$
\State $\alpha \gets 0$
\For{$i = 1$ to $M$}
  \State $\alpha \gets \alpha \cdot n + (c_i - 1)$
  \Comment{Mixed-radix encoding}
\EndFor
\State $\alpha \gets \alpha + 1$
  \Comment{1-indexed}
\State \Return $\alpha$
\end{algorithmic}
\end{algorithm}

\begin{algorithm}[H]
\caption{$\varphi_{31}$: Partitional $\to$ Oscillatory}
\label{alg:phi31}
\begin{algorithmic}[1]
\Require Partition cell index $\alpha \in \{1,\ldots,n^M\}$
\Ensure Oscillatory state $\bm{k} = (k_1, \ldots, k_M) \in
  \{1,\ldots,n\}^M$
\State $r \gets \alpha - 1$
  \Comment{0-indexed}
\For{$i = M$ downto $1$}
  \State $k_i \gets (r \bmod n) + 1$
  \State $r \gets \lfloor r / n \rfloor$
  \Comment{Mixed-radix decoding}
\EndFor
\State \Return $\bm{k}$
\end{algorithmic}
\end{algorithm}

\begin{proposition}[Closure Verification]
$\varphi_{31}(\varphi_{23}(\varphi_{12}(\bm{k}))) = \bm{k}$ for
all $\bm{k} \in \{1,\ldots,n\}^M$.
\end{proposition}

\begin{proof}
$\varphi_{12}$ copies $\bm{k}$ to $\bm{c}$.
$\varphi_{23}$ encodes $\bm{c}$ as a mixed-radix integer $\alpha$.
$\varphi_{31}$ decodes $\alpha$ back to the original tuple.
Mixed-radix encoding/decoding are mutual inverses on
$\{1,\ldots,n\}^M$, so the composition is the identity.
\end{proof}


\section{GPU Memory Budget Breakdown}
\label{app:memory}

Table~\ref{tab:memory} provides a detailed memory budget for the
streaming observation protocol.

\begin{table}[H]
\centering
\caption{Detailed GPU memory budget for the streaming observation
protocol at $256 \times 256$ resolution with float32.}
\label{tab:memory}
\begin{tabular}{lr}
\toprule
\textbf{Component} & \textbf{Size (bytes)} \\
\midrule
Vertex shader (compiled) & 4,096 \\
Observation shader (compiled) & 16,384 \\
Interference shader (compiled) & 8,192 \\
Quality metric shader (compiled) & 12,288 \\
Display shader (compiled) & 8,192 \\
Shader uniform buffers & 1,024 \\
Vertex buffer (fullscreen quad) & 48 \\
\midrule
\textit{Subtotal: shader apparatus} & \textit{50,224}
  $\approx$ 50\,KB \\
\midrule
Query texture $T_q$ ($256^2 \times 4 \times 4$) & 1,048,576 \\
Item texture $T_i$ ($256^2 \times 4 \times 4$) & 1,048,576 \\
Interference texture $T_{\mathrm{int}}$ & 1,048,576 \\
Quality metric texture $T_{\mathrm{qual}}$ & 1,048,576 \\
Gaussian blur intermediate & 1,048,576 \\
\midrule
\textit{Subtotal: textures} & \textit{5,242,880}
  $\approx$ 5\,MB \\
\midrule
Framebuffer objects (2) & 256 \\
Driver overhead (estimated) & 4,194,304 \\
\midrule
\textit{Subtotal: overhead} & \textit{4,194,560}
  $\approx$ 4\,MB \\
\midrule
\textbf{Total} & \textbf{9,487,664}
  $\approx$ \textbf{9.5\,MB} \\
\bottomrule
\end{tabular}
\end{table}

\begin{remark}
The actual GPU memory usage ($\sim$9.5\,MB) is even lower than the
conservative estimate of 13\,MB in Theorem~\ref{thm:memory-bound},
which included a generous overhead margin. In practice, the memory
footprint may range from 8\,MB to 15\,MB depending on the GPU driver
and API (WebGL\,2, Vulkan, Metal).
\end{remark}


\section{S-Entropy Computation Examples}
\label{app:sentropy}

We provide worked examples of S-entropy computation for two domains.

\subsection{Example: Acetone Infrared Spectrum}

Acetone ($\text{CH}_3\text{COCH}_3$) has principal vibrational
modes at approximately:
$\omega_1 = 530\,\text{cm}^{-1}$ (C--C--C bending, $A_1 = 0.3$),
$\omega_2 = 1220\,\text{cm}^{-1}$ (C--C stretching, $A_2 = 0.5$),
$\omega_3 = 1710\,\text{cm}^{-1}$ (C$=$O stretching, $A_3 = 1.0$),
$\omega_4 = 2920\,\text{cm}^{-1}$ (C--H stretching, $A_4 = 0.7$).

Using Eqs.~(\ref{eq:mol-Sk})--(\ref{eq:mol-Se}) with
$\omega_{\max} = 4000\,\text{cm}^{-1}$:

Total intensity: $\sum A_k = 2.5$.
Normalized intensities: $(0.12, 0.20, 0.40, 0.28)$.

$\Sk = -\frac{1}{\ln 4}(0.12 \ln 0.12 + 0.20 \ln 0.20
+ 0.40 \ln 0.40 + 0.28 \ln 0.28) = \frac{1.30}{1.39} \approx 0.94$.

$\St = \frac{0.3 \times 530 + 0.5 \times 1220 + 1.0 \times 1710
+ 0.7 \times 2920}{4000 \times 2.5} = \frac{4524}{10000} \approx 0.45$.

$\bar{\omega} = 4524/2.5 = 1809.6\,\text{cm}^{-1}$.

$\Se = \frac{\sqrt{\sum A_k(\omega_k - 1809.6)^2 / 2.5}}{2000}
\approx \frac{839.1}{2000} \approx 0.42$.

S-entropy triple for acetone: $(\Sk, \St, \Se) \approx (0.94, 0.45, 0.42)$.

Verification: $\Sk + \St + \Se = 1.81$. Note that conservation
(Theorem~\ref{thm:conservation}) applies to the total
\emph{system} entropy, not to the normalized coordinates directly;
the sum of normalized coordinates depends on the normalization
scheme.

\subsection{Example: Human Hemoglobin Protein}

Human hemoglobin ($\alpha$ subunit, 141 residues, PDB: 1HHO) with
contact map $C$ having 847 contacts at 8\,\AA{} cutoff.

Using Eqs.~(\ref{eq:prot-Sk})--(\ref{eq:prot-Se}):

$\Sk = 847 / (141 \times 140/2) = 847/9870 \approx 0.086$.

Amino acid frequencies for the $\alpha$ subunit yield:
$\St = -\frac{1}{\ln 20} \sum_{a=1}^{20} f_a \ln f_a \approx 0.91$.

Contact order: $\sum_{i<j} C_{ij}|i-j|/847 \approx 14.2$, so
$\Se = 14.2/141 \approx 0.10$.

S-entropy triple: $(\Sk, \St, \Se) \approx (0.09, 0.91, 0.10)$.


\section{Complexity Comparison: Formal Derivation}
\label{app:complexity}

\begin{proposition}[Formal Memory Complexity]
\label{prop:formal-memory}
Let $N$ denote the database size and $d$ the feature dimension.

\textbf{Vector database:}
Stores $N$ vectors of dimension $d$:
$M_{\mathrm{VDB}} = N \cdot d \cdot \text{sizeof(float)}
= \Theta(N \cdot d)$.

\textbf{Neural retrieval:}
Stores $N$ vectors plus model parameters $P$:
$M_{\mathrm{NR}} = N \cdot d \cdot \text{sizeof(float)} + P \cdot
\text{sizeof(float)} = \Theta(N \cdot d + P)$.

\textbf{Observation approach:}
Stores shader apparatus $\Mshader$ plus $c$ textures of fixed
resolution $W \times H$:
$M_{\mathrm{obs}} = \Mshader + c \cdot W \cdot H \cdot 4 \cdot
\text{sizeof(float)} = \Theta(1)$.

The observation approach is $\Theta(1)$ in GPU memory because
$\Mshader$, $c$, $W$, and $H$ are all constants independent of $N$.
\end{proposition}

\begin{proof}
$\Mshader$ is determined by the shader source code length and
compiled binary size, both constants. $c = 3$--$5$ (query +
current item + interference + optional quality/display). $W = H = 256$
(or any fixed resolution). None depends on $N$.
\end{proof}

\begin{proposition}[Formal Time Complexity]
\label{prop:formal-time}
For a sequential scan over $N$ items:
$T_{\mathrm{obs}}(N) = N \cdot (T_{\mathrm{upload}} +
T_{\mathrm{render}} + T_{\mathrm{readback}}) = \Theta(N)$,
with constant per-item cost.

For a batched scan with batch size $B$:
$T_{\mathrm{batch}}(N) = \lceil N/B \rceil \cdot T_{\mathrm{batch\_cost}}
= \Theta(N/B) = \Theta(N)$, but with constant reduced by factor $B$.
\end{proposition}


\end{document}